% Memoria técnica de Kyrian World — compilar con: tectonic main.tex (o latexmk -xelatex)
\documentclass[10pt,a4paper,twocolumn]{article}
\usepackage{fontspec}
\setmainfont{texgyrepagella}[Extension=.otf,UprightFont=*-regular,BoldFont=*-bold,ItalicFont=*-italic,BoldItalicFont=*-bolditalic]
\setsansfont{texgyreheros}[Extension=.otf,UprightFont=*-regular,BoldFont=*-bold,ItalicFont=*-italic,BoldItalicFont=*-bolditalic]
\setmonofont{texgyrecursor}[Scale=0.85,Extension=.otf,UprightFont=*-regular,BoldFont=*-bold,ItalicFont=*-italic,BoldItalicFont=*-bolditalic]
\usepackage[spanish,es-tabla,es-noquoting]{babel}
\usepackage[a4paper,margin=18mm,columnsep=7mm]{geometry}
\usepackage{microtype}
\usepackage{booktabs,tabularx,array,multirow,colortbl,xcolor}
\usepackage{graphicx}
\usepackage{tikz}
\usetikzlibrary{arrows.meta,positioning,shapes.geometric,fit,calc,backgrounds}
\usepackage{caption}
\captionsetup{font=small,labelfont=bf}
\usepackage{enumitem}
\setlist{nosep,leftmargin=*}
\usepackage{titlesec}
\usepackage{fancyhdr}
\usepackage[hidelinks]{hyperref}
\usepackage{xurl}
\usepackage[spanish,capitalise,noabbrev]{cleveref}
\crefname{figure}{figura}{figuras}\Crefname{figure}{Figura}{Figuras}
\crefname{table}{tabla}{tablas}\Crefname{table}{Tabla}{Tablas}
\crefname{section}{sección}{secciones}\Crefname{section}{Sección}{Secciones}
\usepackage[numbers,sort&compress]{natbib}
\usepackage{abstract}
\usepackage{stfloats}

\definecolor{kyrian}{HTML}{1F4E5F}
\definecolor{kyrianlight}{HTML}{E8F0F3}
\definecolor{accent}{HTML}{C2571A}
\titleformat{\section}{\large\bfseries\color{kyrian}}{\thesection}{0.6em}{}
\titleformat{\subsection}{\normalsize\bfseries}{\thesubsection}{0.6em}{}
\titlespacing*{\section}{0pt}{1.4ex plus .5ex}{0.8ex}
\titlespacing*{\subsection}{0pt}{1.0ex plus .3ex}{0.4ex}
\pagestyle{fancy}\fancyhf{}
\fancyhead[L]{\small\textsc{Kyrian World} --- Memoria técnica}
\fancyhead[R]{\small\thepage}
\renewcommand{\headrulewidth}{0.3pt}
\newcolumntype{Y}{>{\raggedright\arraybackslash}X}
\newcommand{\code}[1]{\texttt{#1}}
\newcommand{\adr}[1]{ADR~#1}
\setlength{\parskip}{0.2ex}
\setlength{\textfloatsep}{8pt plus 2pt minus 2pt}
\setlength{\dbltextfloatsep}{8pt plus 2pt minus 2pt}
\setlength{\floatsep}{6pt plus 2pt minus 2pt}
\setlength{\intextsep}{6pt plus 2pt minus 2pt}
\setlength{\abovecaptionskip}{4pt}
\setlength{\belowcaptionskip}{0pt}
\setlength{\bibsep}{1pt}

\title{\vspace{-1.2em}\textbf{Kyrian World: un planificador de viajes con IA generativa y recuperación aumentada}\\[0.3em]
\large Proyecto Júpiter --- Trabajo de Fin de Máster\\
\normalsize Máster en Inteligencia Artificial, Cloud Computing y DevOps\\[0.2em]
\normalsize Memoria técnica y de negocio: caso de uso, arquitectura, solución de IA, DevOps, evaluación y trabajo futuro}
\author{Manuel Pérez \and Alexander De Sousa \and Joaquín Castro Salas \and David Baos\\[0.2em] \small Versión 1.1, borrador para revisión --- 23 de septiembre de 2026 --- repositorio: \code{github.com/manupm87/travel-ai-world}}
\date{}

\begin{document}
\twocolumn[
\begin{@twocolumnfalse}
\maketitle
\vspace{-1.5em}
\renewcommand{\abstractname}{Resumen ejecutivo}\begin{abstract}
\noindentKyrian World es un planificador de viajes conversacional, Trabajo de Fin de Máster de un equipo de
cuatro personas apoyado en agentes de IA, que construye itinerarios de una
ciudad a partir de un corpus abierto con licencia limpia (Wikivoyage, Wikipedia, OpenStreetMap)
mediante generación aumentada por recuperación. El MVP está desplegado en AWS en
\code{https://kyrian-world.com}. La memoria plantea: ¿puede un RAG con anclaje
estructural dar itinerarios verificables por poco más de 10~€/mes de coste fijo? En coste, sí; la
calidad del itinerario y la aceptación por usuarios no están medidas. Toda tarjeta procede de un
documento recuperado y los ids no recuperados se descartan; la prosa libre del modelo no está
verificada. El coste fijo estimado es de 12--13~€/mes ($\approx$5~€ en Terraform más un WAF
creado a mano; sin factura real todavía). S3 Vectors se eligió frente a Qdrant por simplicidad
operativa y un \emph{spike} lo confirmó. La viabilidad B2C es marginal; se recomienda el B2B para
oficinas de turismo. La deuda principal es la falta de evaluación extremo a extremo, de alarmas y
de límites de uso.

\end{abstract}
\vspace{1em}
\end{@twocolumnfalse}
]

\section{Introducción}\label{sec:01-introduccion}

Kyrian World es el Trabajo de Fin de Máster (Proyecto Júpiter) de un equipo de cuatro personas,
iniciado el 10-03-2026 con un objetivo concreto: comprobar si un planificador de viajes con IA
generativa puede ofrecer recomendaciones \emph{verificables}, cada una anclada en un documento
real, con un coste de operación de unos pocos euros al mes. Hasta el 23-09-2026 acumula 324
\emph{commits} y 181 PR fusionadas, está desplegado en AWS (\code{kyrian-world.com}) y se ha
desarrollado en buena parte con agentes de programación guiados por un proceso escrito
(\emph{brief}, implementación, revisión y corrección), con cada cambio trazable hasta un issue y,
cuando procede, hasta una ADR.

\paragraph{Objetivos y alcance.} Explicar las decisiones de producto, arquitectura y plataforma
con su contexto y alternativas; valorar la viabilidad en su mercado; y cerrar con un análisis
crítico y propuestas priorizadas. Se describe el sistema en el \emph{commit} \code{e7916c8}.
Quedan fuera la alternativa en GCP, que no se despliega, y cualquier medida con usuarios reales,
que no existen todavía.

\paragraph{Metodología.} Los hechos proceden del repositorio: código, 23 ADR, \emph{runbooks}, el
documento del \emph{spike} de vectores y el historial de \code{git} y de GitHub. Cada cifra indica
si es \emph{medida} (el \emph{spike}; los tests, ejecutados con \code{just test} el 23-09-2026),
\emph{estimada} (costes de los ADR, economía unitaria) o \emph{contada} (líneas, PR, documentos
del corpus); lo no medido se declara. El estudio de mercado usa fuentes públicas citadas y marca
sus proyecciones como hipótesis. La \cref{tab:01-guion} relaciona cada apartado y requisito
técnico del guion con la sección que lo cubre.

\begin{table}[tb]
\centering\footnotesize
\caption{Correspondencia con el guion del Proyecto Júpiter.}\label{tab:01-guion}
\begin{tabularx}{\columnwidth}{@{}Yll@{}}
\toprule
Apartado o requisito del guion & Dónde & Estado \\
\midrule
a) Introducción y resumen ejecutivo & Resumen, \cref{sec:01-introduccion} & cubierto \\
b) Caso de negocio y mercado & \cref{sec:02-producto,sec:03-mercado} & cubierto \\
c) Contribución individual & \cref{sec:03b-equipo} & cubierto \\
d) Arquitectura técnica & \cref{sec:04-arquitectura-software,sec:06-nube} & cubierto \\
e) Diseño de la solución de IA & \cref{sec:05-arquitectura-ia} & cubierto \\
f) Integración con DevOps & \cref{sec:07-entorno-desarrollo,sec:08-cicd-proceso} & cubierto \\
g) Uso de modelos preentrenados & \cref{sec:05-arquitectura-ia} & cubierto \\
h) Evaluación de la solución & \cref{sec:09b-evaluacion} & parcial \\
i) Conclusiones y trabajo futuro & \cref{sec:10-analisis-critico,sec:11-conclusiones} & cubierto \\
\midrule
IA generativa & \cref{sec:05-arquitectura-ia} & cubierto \\
Base de datos vectorial & \cref{sec:05-arquitectura-ia} & cubierto \\
MVP funcional & \cref{sec:02-producto,sec:06-nube} & cubierto \\
API de modelo fundacional & \cref{sec:05-arquitectura-ia} & cubierto \\
Pipeline CI/CD & \cref{sec:08-cicd-proceso} & cubierto \\
Monitorización y trazabilidad & \cref{sec:09-sre} & parcial \\
\bottomrule
\end{tabularx}
\end{table}

\paragraph{Estructura.} Las \cref{sec:02-producto,sec:03-mercado} presentan el producto, el
mercado y la viabilidad, y la \cref{sec:03b-equipo}, el equipo y la contribución de cada miembro.
Las \cref{sec:04-arquitectura-software,sec:05-arquitectura-ia,sec:06-nube} describen las
arquitecturas de software, de IA y de nube; las
\cref{sec:07-entorno-desarrollo,sec:08-cicd-proceso,sec:09-sre} tratan el entorno de desarrollo,
la entrega continua con agentes y la fiabilidad y los costes, y la \cref{sec:09b-evaluacion}, la
evaluación. La \cref{sec:10-analisis-critico} reúne la crítica y las propuestas, la
\cref{sec:11-conclusiones} concluye y el anexo (\cref{sec:12-anexo}) recoge un glosario y la tabla
de ADR.

\section{El producto}\label{sec:02-producto}

Kyrian World (\code{kyrian-world.com}) es un planificador de viajes conversacional que construye
en tiempo real un itinerario día a día que se puede guardar y reabrir. El usuario objetivo, inferido del alcance (no hay estudio de usuarios), es el viajero
individual en \emph{una sola ciudad} (\adr{0019}). A diferencia de un chat genérico, cada
recomendación es una tarjeta anclada en un documento del corpus (Wikivoyage, Wikipedia,
OpenStreetMap), con foto y atribución (\cref{sec:05-arquitectura-ia}).

En el recorrido (\cref{fig:02-flujo,tab:02-rutas}), el login (Google en local, Cognito con
PKCE en producción) conserva el destino \code{/plan/?q=} y \code{?trip=<uuid>} reconstruye el
borrador. Sin \code{ai\_api}, un \emph{modo demo} avisado reproduce una sesión de Budapest; la
consola \code{admin} es de solo lectura (\adr{0024}).

\begin{figure}[tb]\centering
\begin{tikzpicture}[
  n/.style={draw, rounded corners, fill=kyrianlight, font=\small, align=center,
            minimum height=8mm, text width=2.9cm},
  s/.style={n, fill=white, dashed, text width=2.5cm},
  a/.style={-{Stealth}}, l/.style={font=\footnotesize, align=center},
  node distance=5mm and 10mm]
\node[n] (land) {Portada \code{/}\\ campo de pregunta};
\node[n, below=of land] (login) {Inicio de sesión\\ Google / Cognito};
\node[n, below=of login] (home) {Mis viajes\\ \code{/dashboard/}};
\node[n, below=of home] (plan) {Planner \code{/plan/}\\ chat $\cdot$ viaje $\cdot$ mapa};
\node[n, below=of plan] (trip) {Viaje guardado\\ \code{?trip=<uuid>}};
\node[s, right=of login] (admin) {Consola admin\\ \code{/admin/} (rol)};
\node[s, right=of plan] (demo) {Modo demo\\ sesión grabada};
\node[s, right=of trip] (lock) {En curso o pasado:\\ solo lectura};
\draw[a] (land) -- node[l, left]{sin sesión} (login);
\draw[a] (login) -- (home);
\draw[a] (home) -- node[l, left]{nuevo / abrir} (plan);
\draw[a] (plan) -- node[l, left]{guardar} (trip);
\draw[a] (login) -- node[l, above]{admin} (admin);
\draw[a, dashed] (demo) -- node[l, above]{sin API} (plan);
\draw[a] (trip) -- node[l, above]{fase} (lock);
\draw[a] (land.west) -- ++(-3mm,0) |- node[l, pos=0.25, left]{con\\sesión} (plan.west);
\end{tikzpicture}
\caption{Recorrido del usuario. En trazo discontinuo, superficies laterales al flujo principal.}
\label{fig:02-flujo}
\end{figure}

\begin{table}[tb]\centering\footnotesize
\begin{tabularx}{\columnwidth}{@{}lY@{}}\toprule
Ruta & Función \\\midrule
\code{/} & Portada: un único campo de pregunta \\
\code{/auth/callback/} & Retorno del login gestionado de Cognito \\
\code{/dashboard/} & Mis viajes: campo de pregunta y viajes por fase \\
\code{/plan/} & Planner (\code{?q=} o \code{?trip=}): chat, viaje y mapa \\
\code{/trip/?id=} & Redirección heredada a \code{/plan/?trip=} \\
\code{/admin/} & Consola: KPIs por día, modelo y ciudad \\
\code{/admin/…} & \code{turns/}+\code{turn/} (inspector), \code{trips/}+\code{trip/}, \code{users/} (rol y \code{subject}) \\
otra & 404 exportado como \code{404.html} \\\bottomrule
\end{tabularx}
\caption{Páginas del export estático. Los identificadores van en la \emph{query string} porque
un export estático no conoce en compilación los ids de los usuarios (\adr{0011}).}
\label{tab:02-rutas}
\end{table}

\subsection{El planner}

En escritorio, tres columnas (chat $\approx$30\,\%, viaje $\approx$40\,\%, mapa $\approx$30\,\%);
en móvil, pestañas. Un \emph{brief} (ciudad, fechas, origen, presupuesto, intereses), completado con
preguntas cortas y chips, da paso a carruseles de tarjetas: barrios, hoteles por nivel de precio (\code{€} a
\code{€€€}; nunca un número) y actividades por franja (mañana, tarde, atardecer, noche), con
foto acreditada, enlace a la fuente y ``Cambiar'' para ver alternativas. Los lugares
citados en prosa llegan ``sin colocar'' y el viajero elige día y franja: el sistema nunca mueve
su mañana sin pedirlo (\adr{0018}).

El panel central muestra la ruta (enlace a Google Flights, sin precio), la estancia, una tira
de días con previsión (Open-Meteo~\cite{fe:openmeteo} a corto plazo, normales climáticas del
corpus más allá) y una vista general de la ciudad. El mapa (MapLibre GL~\cite{fe:maplibre}, teselas de
OpenFreeMap~\cite{fe:openfreemap}) numera las paradas del día unidas por líneas rectas, sin
cálculo de rutas.

Guardar crea un \code{Trip} en \code{core\_api}. \emph{Contexto}: el modelo
previo era multidestino con estado almacenado. \emph{Decisión} (\adr{0019}): una ciudad por viaje
y fase (próximo, en curso, pasado) derivada de las fechas en cada respuesta; un viaje en curso o
pasado es de solo lectura (409 \code{TRIP\_LOCKED}, sin campo de texto). \emph{Consecuencia}: la
fase no puede desincronizarse, pero la migración borró los viajes multidestino y la fase se
calcula en UTC. \emph{Valoración}: recorte legítimo para un producto mínimo, pero excluye interrail y rutas
multiciudad. \adr{0020} creó ``Mis viajes'' días después porque el login dejaba un planner
vacío: la entrada no se validó antes.

\subsection{Identidad, móvil y accesibilidad}

La identidad (\code{docs/design/kyrian-world.md}) parte de que ``el campo es lo memorable'':
sin rejilla de funciones ni testimonios, solo un fondo de ``aurora'' (\code{aria-hidden}), tema
oscuro, variables CSS, Outfit y Plus Jakarta Sans. La consola, a propósito, es densa y usa
gráficos SVG hechos a mano.

Inglés y español cambian sin recarga; tipos y un test exigen las mismas claves. En móvil (TRA-187), \code{100dvh} evita que la barra del navegador tape el campo, cada panel
desplaza por separado y los controles miden 16\,px en punteros táctiles (sin zoom de iOS),
comprobado por E2E a 390$\times$844.
Según WCAG~2.2~\cite{fe:wcag22}: contraste $\geq$4,5:1 calculado token a token
(el dorado del tema claro bajó a \code{\#8A5400}, 5,8:1, tras descartar dos tonos), un único
contrato de diálogo (foco atrapado, Escape lo devuelve), \code{prefers-reduced-motion} global y
tablas \code{sr-only} tras los gráficos de la consola. Sin auditoría automatizada (axe,
Lighthouse) en CI, el cumplimiento está declarado, no medido.

\subsection{Fotos: mejor ninguna que la equivocada}

\emph{Contexto}: tiene foto el 80--84\,\% de los documentos de monumentos y visitas, pero solo el
0--2\,\% de restaurantes y el 6--12\,\% de hoteles, y el recurso previo mostraba la foto de
\emph{otro} local. Google Places Photos y TripAdvisor se descartaron por sus términos.
\emph{Decisión} (\adr{0021}): cascada imagen del corpus $\to$ Wikimedia Commons $\to$
\code{og:image} del sitio del propio local (no persistida, con listas de bloqueo) $\to$ foto de un monumento solo para barrios $\to$ marcador neutro. Para hoteles,
\adr{0022} resuelve la foto al construir el corpus (sitio propio, Facebook, Commons por nombre,
mayor imagen de la portada) y \emph{un hotel sin foto sale del corpus}: con foto localizada,
Budapest pasó de 36 a 206 de 415, Berlín 338/649 y Madrid 279/527 (antes del filtrado; la \cref{tab:05-corpus} da los que quedan tras él). \emph{Valoración}: confianza antes que inventario, a costa de casi la mitad de los hoteles,
de enlaces directos que pueden romperse hasta la siguiente reconstrucción y de una capa curada a
mano (cadenas) sin verificación automática.

\paragraph{Stack.} Export estático de Next.js~16~\cite{fe:nextjs} en S3 y CloudFront
(\cref{sec:06-nube}). El lint impone que \code{fetch} y los tipos generados desde OpenAPI solo
vivan en \code{src/services/}, donde un mapeador los convierte en modelos de vista (\adr{0006});
el estado del planner es un reductor puro probado sin React (\cref{tab:02-stack}).

\begin{table}[tb]\centering\footnotesize
\begin{tabularx}{\columnwidth}{@{}lY@{}}\toprule
Elemento & Valor \\\midrule
Framework & Next.js 16.3, React 19.3, TypeScript 5.9 estricto \\
Estilos, mapa & Tailwind v4 sin \code{tailwind.config.js}; MapLibre GL 6.10 \\
Pruebas & Vitest 5, Testing Library, Playwright 1.63 (3 configuraciones) \\
\code{.ts}/\code{.tsx} & 259 ficheros (83 de test); 34\,623 líneas (22\,333 sin tests) \\
Tipos generados & 2 ficheros aparte, 4\,664 líneas (\code{core-api.ts}, \code{ai-api.ts}) \\
Tests unitarios & 747 casos en 83 ficheros, todos pasan (Vitest, 23-09-2026) \\
Specs E2E & 9 ficheros, 1\,917 líneas (no ejecutados para esta memoria) \\\bottomrule
\end{tabularx}
\caption{Stack y tamaño del frontend (\code{package.json}, recuento y ejecución en \code{src/frontend}).}
\label{tab:02-stack}
\end{table}

\paragraph{Valoración.} El producto resuelve un caso acotado (una ciudad, un viajero, hasta
una semana). Las tarjetas no se reordenan, no hay horarios ni notas, y guardar reescribe una
instantánea completa (TRA-146). No hay colaboración, reservas, precios reales, PDF ni aplicación
instalable. El mapa depende de un servicio gratuito sin SLA y el emparejamiento de lugares en la prosa es heurístico. Faltan mejoras obvias: enseñar en la
tarjeta el fragmento recuperado, que es el diferenciador (hoy solo el enlace a la fuente, con la
licencia en un \emph{tooltip}); tarjetas en el idioma del usuario (hoy en inglés); tiempos de
desplazamiento entre paradas, y una valoración por tarjeta ligada a las trazas
(\cref{sec:10-analisis-critico}).

\section{Estudio de mercado y viabilidad}\label{sec:03-mercado}

\subsection{Mercado y tendencias}

El turismo internacional cerró 2025 con 1\,520 millones de llegadas (+4\,\%); Europa recibió
793 millones \cite{mk:untourism2025}.
España batió su récord con 96,8 millones de turistas internacionales y 134\,712~M€ de gasto
\cite{mk:egatur2025}. Las OTA facturaron 94\,000~M\$ en 2024 (+7\,\% anual) y ven
en la IA conversacional una amenaza \cite{mk:skiftota}. Los ingresos se generan al \emph{reservar};
\emph{planificar}, el eslabón que cubre Kyrian World, apenas se monetiza por sí mismo.

En España el 40\,\% de los adultos ya ha usado un asistente de IA para preparar un viaje (el
73\,\% entre los 18 y los 24 años) \cite{mk:yougov2026}; Europa va por detrás (Reino Unido
22\,\%, Francia 19\,\%, Alemania 15\,\%) \cite{mk:phocuseurope2025}. Falta confianza: solo el 17--20\,\% de los viajeros europeos confía plenamente en una IA
\cite{mk:phocustrust}, y el 49\,\% de los españoles que no la usan teme recomendaciones
inexactas o desactualizadas \cite{mk:yougov2026}. En una encuesta patrocinada por Mindtrip a
103 organizaciones de gestión de destinos (DMO) de EE.\,UU., solo el 22\,\% ofrece un
planificador con IA en su web \cite{mk:dmoreadiness}.

\subsection{Competencia}

La \cref{tab:03-competidores} enseña que los generadores sin datos propios acaban comprados y
que Expedia retiró Romie por responder sin inventario ni precios reales \cite{mk:romie}.

\begin{table*}[tb]\centering\footnotesize
\begin{tabularx}{\textwidth}{@{}p{1.9cm}p{2.4cm}YYY@{}}\toprule
Nombre & Tipo / propietario & Negocio & Diferenciador & Debilidad o estado \\\midrule
Layla & \emph{Startup} ($\approx$5~M€) & \emph{Freemium} (9,99~\$/mes), afiliación \cite{mk:laylapricing} & Chat y vídeo; precios en vivo & Comprada por Expedia (07-2026); mantiene su marca \cite{mk:laylaexpedia} \\
Roam Around & \emph{Startup} (1~M\$) & Gratuito & Sencillez; auge tras ChatGPT & Absorbida por Layla (02-2024); sin datos propios \\
Mindtrip & \emph{Startup} (19~M\$) & Gratis; comisión; B2B & 11~M de POI; >35 destinos clientes \cite{mk:mindtrip} & Fuerte en EE.\,UU.; poco español \\
GuideGeek & Matador Network & Marca blanca para DMO & >70 DMO; WhatsApp, Instagram \cite{mk:guidegeek} & Depende de OpenAI; generalista \\
Wonderplan & \emph{Startup} & Gratis, sin registro & Rapidez; PDF & Casi insensible al perfil \\
Wanderlog & \emph{Startup} (YC) & \emph{Freemium} & Colaboración y mapas & IA secundaria \\
ChatGPT & OpenAI & Suscripción; apps de Booking y Expedia (10-2025) \cite{mk:chatgptapps} & Distribución masiva & Sin anclaje; reserva agéntica experimental \\
Google AI Mode & Alphabet & Publicidad, comisión & Itinerarios (11-2025); datos de Maps \cite{mk:googleaimode} & Solo EE.\,UU.; Google Trips cerró en 2019 \\
Expedia Romie & Expedia Group & Comisión OTA & Planificación y soporte (2024) & Retirado: sin inventario ni precios reales \cite{mk:romie} \\
Booking.com & Booking Holdings & Comisión OTA & Chat sobre inventario propio (2023) & Ligado a su inventario; app en ChatGPT \\\bottomrule
\end{tabularx}
\caption{Competidores y comparables (septiembre de 2026). Las cifras sin cita proceden del
estudio de mercado del proyecto (reseñas y prensa sectorial).}
\label{tab:03-competidores}
\end{table*}

\textbf{Posicionamiento.} Kyrian World compite en \emph{verificabilidad}, no en distribución ni
inventario: cada tarjeta remite, con enlace a la fuente, a un documento de un corpus con licencia
limpia (\cref{sec:05-arquitectura-ia}). Es un anclaje documental (existencia, descripción, fuente),
no de disponibilidad ni precio: cubrir la carencia de Romie exige integrar inventario (horizonte
de 12 meses). Asume una ciudad por viaje (\adr{0019}) y pocas ciudades, aunque en castellano. Su
ventaja es el coste fijo: $\approx$5~€/mes en recursos de Terraform (\adr{0023}) más
$\approx$8~\$/mes de un WAF creado a mano, unos 12--13~€/mes (estimación; la factura de septiembre
debe sustituirla), frente a rondas de millones.

\subsection{Monetización}

La \textbf{afiliación} paga por reserva completada (\cref{tab:03-afiliacion}), pero el modelo por sesión de Booking.com solo computa la reserva hecha en el
navegador inicial \cite{mk:bookingcj}. La \textbf{suscripción} toma como referencia Layla
Premium, a 9,99~\$/mes o 49,99~\$/año \cite{mk:laylapricing}; se supone un plan de 39~€/año con
una conversión del 1--2\,\% (\emph{hipótesis}: el contenido del plan \emph{premium} está sin definir). La vía con más recorrido es el \textbf{B2B o
marca blanca para DMO}, con demanda demostrada por GuideGeek \cite{mk:guidegeek} y Mindtrip
\cite{mk:mindtrip}: se vendería el \emph{pipeline} \code{/add-city}, a 6\,000--15\,000~€ por
destino y año (\emph{hipótesis}).

\begin{table}[tb]\centering\footnotesize
\begin{tabularx}{\columnwidth}{@{}lY@{}}\toprule
Programa & Comisión de referencia \\\midrule
Booking.com & 4\,\% de estancias completadas; por sesión \cite{mk:bookingcj} \\
Expedia Group & Hoteles 3--4\,\%; actividades 4\,\% \cite{mk:expediaaff} \\
GetYourGuide / Viator & $\approx$8\,\%; \emph{cookie} de 30 días \cite{mk:viator} \\
Kiwi.com & 3\,\% del billete \cite{mk:kiwi} \\
Omio & 2--8\,\% según mercado y volumen \cite{mk:omio} \\\bottomrule
\end{tabularx}
\caption{Comisiones de afiliación de referencia (consultadas el 23-09-2026).}
\label{tab:03-afiliacion}
\end{table}

\subsection{Economía unitaria}

Con unas 12 llamadas al modelo por viaje, $\approx$50\,000 \emph{tokens} de entrada y
$\approx$8\,000 de salida (\emph{estimación}), el coste variable por viaje es de 0,09~\$ con
Claude Haiku 4.5, el modelo de producción; 0,18--0,30~\$ con Sonnet (el \emph{endpoint} de la UE
es un 10\,\% más caro) \cite{mk:anthropicpricing}; 0,07~\$ con Nova Pro \cite{mk:bedrockpricing};
vectores y \emph{embeddings}, $<$0,001~\$. El escenario usa 0,20~€/viaje para dejar margen a
reintentos, llamadas de chat y un posible modelo mayor.

El ingreso por afiliación es de 0,37~€ por viaje en el escenario base (0,15--0,76~€), con tasas
de clic y conversión \emph{supuestas}: el margen queda en $\approx$0,17~€ por viaje, y en el
escenario pesimista cada viaje pierde dinero. En B2C, el LTV ronda 1,8~€ (\emph{hipótesis}:
1,8 viajes/año $\times$ 0,17~€ más 1,5\,\% $\times$ 39~€ de \emph{premium}, $\approx$0,9~€/año
durante dos años). Como CAC se usa de \emph{proxy} el coste por instalación de una app móvil de
viajes, 4,20~\$ \cite{mk:cac}: LTV/CAC $\approx$0,5, así que el B2C depende de canales orgánicos. En B2B, un
contrato de 8\,000~€/año durante tres años frente a 2\,000--4\,000~€ de coste de venta pública da
un LTV/CAC de 6--12 (\emph{hipótesis}).

\begin{table}[tb]\centering\footnotesize
\begin{tabularx}{\columnwidth}{@{}Yrrr@{}}\toprule
Concepto (€) & Año 1 & Año 2 & Año 3 \\\midrule
Viajes planificados & 5\,200 & 40\,000 & 162\,000 \\
Afiliación (0,37~€/viaje) & 1\,924 & 14\,800 & 59\,940 \\
\emph{Premium} (1--2\,\% $\times$ 39~€) & 1\,560 & 14\,625 & 70\,200 \\
B2B (contratos) & 3\,000 (1) & 32\,000 (4) & 108\,000 (12) \\
\textbf{Ingresos} & \textbf{6\,484} & \textbf{61\,425} & \textbf{238\,140} \\
\quad peso del B2B & 46\,\% & 52\,\% & 45\,\% \\\midrule
Coste variable (0,20~€/viaje) & 1\,040 & 8\,000 & 32\,400 \\
Infra., datos y legal & 1\,980 & 7\,440 & 19\,800 \\
Marketing y venta B2B & 3\,000 & 28\,000 & 95\,000 \\
\textbf{Resultado antes de salarios} & \textbf{+464} & \textbf{+17\,985} & \textbf{+90\,940} \\
Salarios & 0 & 35\,000 & 90\,000 \\
\textbf{Resultado neto} & \textbf{+464} & \textbf{$-$17\,015} & \textbf{+940} \\\bottomrule
\end{tabularx}
\caption{Escenario base a tres años (10-2026 a 09-2029). Hipótesis: 4\,000 / 25\,000 / 90\,000
usuarios activos con 1,3 / 1,6 / 1,8 viajes al año; B2B a 3\,000~€ (piloto), 8\,000~€ y
9\,000~€ por contrato; fundador asalariado desde el año 2 y una segunda persona en el año 3.}
\label{tab:03-escenario}
\end{table}

En la \cref{tab:03-escenario} el año~1 cuadra solo porque no imputa salario: con un salario
mínimo imputado daría $\approx-$15\,000~€ (\emph{hipótesis}). Cubrir 35\,000~€ de salario y 5\,000~€ de fijos
exigiría, solo con afiliación, unos 235\,000 viajes al año, irreal sin adquisición de pago; con
B2B bastan cinco contratos de $\approx$8\,000~€.

\subsection{Marco legal}

Entre corchetes, el estado de cada obligación. El \textbf{RGPD} exige base jurídica,
minimización, plazos de conservación de las trazas y cuidado con el texto libre, que puede revelar
categorías especiales \cite{mk:anthropicpricing}; [\emph{pendiente}: el login menciona una política de privacidad y unos
términos cuyas páginas no existen]. El \textbf{Reglamento de IA} (UE) 2024/1689 no incluye un
planificador de viajes en el Anexo~III: le afecta el art.~50 (desde el 02-08-2026): avisar de que
se habla con una IA y marcar de forma legible por máquina el contenido sintético \cite{mk:aiact50} [\emph{parcial}: el planner se
titula ``Dile a la IA adónde quieres ir'', sin aviso explícito ni marcado]. La \textbf{Directiva de
viajes combinados} no aplica mientras Kyrian World solo redirija. Las \textbf{licencias} imponen
atribución: CC BY-SA 4.0 (Wikipedia, Wikivoyage) obliga a compartir igual las adaptaciones
publicadas, como el corpus JSONL \cite{ai:wikivoyagelicense}; ODbL exige atribuir a OpenStreetMap \cite{ai:osmcopyright}
[\emph{parcial}: la tarjeta enlaza la fuente, pero la licencia solo aparece en un \emph{tooltip};
la introducción de la ciudad sí cita ``Wikivoyage (CC BY-SA 4.0)'']. La API gratuita de
\textbf{Open-Meteo} es solo no comercial \cite{mk:openmeteo} [\emph{pendiente}: monetizar exige su
plan comercial u otra fuente].

\textbf{Vender a una DMO pública española.} La licencia hipotética de 6\,000--15\,000~€ encaja en
un contrato menor de servicios (art.~118 LCSP, $<$15\,000~€) \cite{fx:lcsp}. Exigiría además
conformidad con el Esquema Nacional de Seguridad \cite{fx:ens}, sin evaluar; accesibilidad según el
RD~1112/2018 y la EN~301~549 \cite{fx:rd1112,fx:en301549}, declarada pero no auditada
(\cref{sec:02-producto}); alojamiento en la UE, que la región \code{eu-west-1} cubre para la API
pero sin acreditación formal; contenido en español, cuando el 98,8\,\% del corpus está en inglés, y
cobertura de ciudades medianas, donde Bolonia tiene solo 1\,597 documentos. Hoy ninguno está
acreditado.

\begin{table}[tb]\centering\footnotesize
\begin{tabularx}{\columnwidth}{@{}YY@{}}\toprule
\textbf{Fortalezas} & \textbf{Debilidades} \\\midrule
\begin{itemize}[leftmargin=*,label=--,nosep]
\item Anclaje citado, licencia limpia
\item 12--13~€/mes de coste fijo
\item Alta de ciudad reproducible
\end{itemize} &
\begin{itemize}[leftmargin=*,label=--,nosep]
\item Pocas ciudades; sin reservas
\item Sin métricas ni ingresos
\item \emph{Bus factor} 1
\end{itemize} \\\midrule
\textbf{Oportunidades} & \textbf{Amenazas} \\\midrule
\begin{itemize}[leftmargin=*,label=--,nosep]
\item 78\,\% de DMO sin IA (EE.\,UU.)
\item Red DTI: 440 destinos
\item Art.~50; LLM más barato
\end{itemize} &
\begin{itemize}[leftmargin=*,label=--,nosep]
\item Google y ChatGPT gratis
\item Expedia compra Layla
\item Comisiones y datos abiertos
\end{itemize} \\\bottomrule
\end{tabularx}
\caption{DAFO de Kyrian World (septiembre de 2026).}
\label{tab:03-dafo}
\end{table}

\subsection{Conclusión de viabilidad}

Como B2C financiado con afiliación (\cref{tab:03-dafo}), la viabilidad es marginal: el ingreso
por viaje es del orden del coste del modelo, el LTV no cubre el CAC y las plataformas regalan la planificación. El nicho recomendado
es el B2B para oficinas de turismo de ciudades europeas medianas y los 440 destinos de la Red de
Destinos Turísticos Inteligentes \cite{mk:segittur}, con el viajero hispanohablante como segmento
secundario. El diferenciador a defender es la \emph{IA verificable con fuentes}, en línea con la
transparencia del art.~50; la carencia de Romie \cite{mk:romie}, inventario y precios reales, sigue
sin cubrir.

\section{Equipo y contribución individual}\label{sec:03b-equipo}

El equipo lo forman cuatro personas. La \cref{tab:03b-equipo} resume, \emph{según git y GitHub}
(PR fusionadas en \code{main} hasta \code{e7916c8}), qué hizo cada una; el rol se deduce de las
rutas del repositorio que tocó y de los títulos de sus PR, no de un reparto formal previo.

\begin{table}[tb]
\centering\footnotesize
\caption{Contribución por miembro según las PR fusionadas (181 en total, 3 de ellas de Dependabot).}\label{tab:03b-equipo}
\begin{tabularx}{\columnwidth}{@{}>{\raggedright\arraybackslash}p{1.55cm}Y@{}}
\toprule
Miembro (PR) & Rol principal, áreas del repositorio e hitos (n.º de PR) \\
\midrule
Manuel Pérez (145) & Producto de extremo a extremo y proceso de entrega con agentes.
\code{src/frontend}, \code{core\_api}, planificador de \code{ai\_api}, \code{city\_corpus},
migración a DynamoDB, consola de administración (\#193--\#196), \code{infra/aws} v3, CI/CD,
\code{docs} y ADR. Por prefijo del título: $\approx$27 infraestructura/CI, 22 \emph{frontend},
21 \code{core\_api}, 20 documentación, 17 corpus, 17 \code{ai\_api}, 8 administración, 13 otras. \\
\addlinespace
Alexander De Sousa (19) & Prototipo del planificador (\#17--\#20); \emph{scraper} de Google
Places y Wikipedia (\#26--\#31, \#68); \emph{devcontainer} inicial (\#35, \#36); Terraform
inicial de AWS y GCP (\#37); \emph{frontend} HTTPS con S3, CloudFront, Route~53 y ACM (\#83);
corpus de Madrid (\#172). \\
\addlinespace
Joaquín Castro Salas (14) & Andamiaje del \emph{backend} y esquema de datos (\#21--\#25); OAuth de
Google, chat, CI/CD (\#32); protección del chat y CI del \emph{backend} (\#40, \#41); Bedrock en
\code{ai\_api} (\#109); hilos de chat (\#111, \#115); infraestructura de S3 Vectors (\#116); RAG
del chat con S3 Vectors y Titan (\#124); corpus de Berlín (\#171). \\
\addlinespace
David Baos (0; 1 rama) & Ciudad de Miami a través del \emph{pipeline} de \code{city\_corpus}
(TRA-204, rama \code{tra-204-miami-through-the-pipeline}, 21-09-2026): configuración
\code{cities/miami.toml}, corpus de 1\,249 documentos, manifiesto e informe de \emph{readiness};
sin PR fusionada en \code{main} en la fecha de corte. \\
\bottomrule
\end{tabularx}
\end{table}

\paragraph{Cómo se registró la contribución.} Todo cambio sigue la misma cadena: issue en Linear
(TRA-$n$) $\rightarrow$ rama \code{<tipo>/TRA-<n>-<slug>} $\rightarrow$ PR con autor en GitHub
$\rightarrow$ fusión. El código generado con agentes lleva en el \emph{commit} la línea
\code{Co-Authored-By: Claude} (151 de los 324 \emph{commits} de \code{main}); aun así, cada PR la
abrió, revisó y fusionó una persona, que es su autora a efectos de esta tabla
(\cref{sec:08-cicd-proceso}).

\paragraph{Limitaciones de la evidencia.} El reparto es desigual (145/19/14/0 PR fusionadas) y el
recuento de PR o de líneas no mide esfuerzo: no captura el trabajo no versionado (reuniones,
curación manual de datos, pruebas en producción) ni el peso de cada PR, y deja fuera el trabajo
que aún no ha llegado a \code{main}, como la ciudad de Miami.

\section{Arquitectura de software}\label{sec:04-arquitectura-software}

Un frontend Next.js exportado como sitio estático y dos servicios FastAPI \cite{sw:fastapi} que
solo comparten la forma de verificar el token (\cref{fig:04-contenedores}). El pipeline de IA se
trata en \cref{sec:05-arquitectura-ia} y el despliegue en \cref{sec:06-nube}.

\begin{figure*}[tb]
\centering
\begin{tikzpicture}[
  box/.style={draw, rounded corners, fill=kyrianlight, font=\small, align=center, minimum height=8mm, text width=#1},
  box/.default=2.9cm,
  store/.style={draw, rounded corners, fill=white, font=\small, align=center, minimum height=8mm, text width=3.3cm, line width=0.8pt, draw=kyrian},
  arr/.style={-{Stealth}, thick},
  lbl/.style={font=\footnotesize, align=center, fill=white, inner sep=1pt}]
\node[box=2.2cm] (nav) at (0,0) {Navegador\\{\footnotesize Next.js estático}};
\node[box=2.7cm] (px) at (4.3,0) {CloudFront (AWS)\\ nginx (Compose)};
\node[box=3.0cm, fill=kyrian!18] (core) at (8.3,1.35) {\code{core\_api}\\{\footnotesize capas · auth · viajes · chat}};
\node[box=3.0cm, fill=accent!15] (ai) at (8.3,-1.35) {\code{ai\_api}\\{\footnotesize puertos y adaptadores · SSE}};
\node[box=2.7cm] (idp) at (4.3,2.45) {Cognito + Google IdP\\{\footnotesize local: Google \code{tokeninfo}}};
\node[store] (ddb) at (13.5,1.95) {DynamoDB \code{<prefix>-core}\\{\footnotesize tabla única (\adr{0023})}};
\node[store] (llm) at (13.5,0.35) {LLM\\{\footnotesize Bedrock (nube) · NVIDIA (local)}};
\node[store] (vec) at (13.5,-1.05) {S3 Vectors\\{\footnotesize + embeddings Titan V2}};
\node[store] (itx) at (13.5,-2.45) {DynamoDB \code{<prefix>-interactions}\\{\footnotesize trazas (\adr{0024})}};
\draw[arr] (nav) -- node[lbl, above=1pt] {\code{/} y \code{/api/*}} (px);
\draw[arr] (px.north east) -- node[lbl, above left=-1pt] {\code{/api/*}} (core.west);
\draw[arr] (px.south east) -- node[lbl, below left=-1pt] {\code{/api/v1/ai/*}} (ai.west);
\draw[arr] (nav.north) |- node[lbl, pos=0.75, above] {código + PKCE} (idp.west);
\draw[arr, dashed] (core.north) |- node[lbl, pos=0.75, above] {solo modo local} (idp.east);
\draw[arr] (core.east) -- (ddb.west);
\draw[arr] (ai.east) -- (llm.west);
\draw[arr] (ai.east) -- node[lbl, pos=0.55] {IAM} (vec.west);
\draw[arr] (ai.east) -- (itx.west);
\draw[arr, very thick, draw=accent] (ai.north) -- node[lbl, right=2pt, text width=3.6cm, align=left] {HTTP con el \emph{bearer} del\\ usuario; sentido único\\ \code{ai\_api} $\rightarrow$ \code{core\_api}} (core.south);
\end{tikzpicture}
\caption{Contenedores y flujos. \code{core\_api} funciona con \code{ai\_api} caído; ambos
verifican el JWT por su cuenta (HS256 local o RS256 con el JWKS de Cognito) y cada almacén tiene
un único servicio propietario.}
\label{fig:04-contenedores}
\end{figure*}

\subsection{Dos servicios, una frontera y un núcleo compartido}

\adr{0001} dividió el FastAPI inicial (unas 3\,150 líneas) porque el RAG traía dependencias
pesadas, \emph{streaming} largo y otro perfil de escalado, y el único acoplamiento era
\code{get\_current\_user} leyendo la base de datos. La frontera se verifica:
\code{ai\_api} y \code{core\_api} no se importan entre sí, y
\code{tests/test\_import\_boundaries.py} comprueba además que solo \code{infrastructure/dynamo/}
importa boto3. Lo compartido vive en \code{travel\_common}: \code{Principal}, ajustes, errores de
dominio, verificación de tokens, \emph{app factory} y \emph{fakes}. Todo forma un \emph{workspace} uv con un único \code{uv.lock}, ruff con reglas ampliadas
(\code{E4/E7/E9/F}, \code{I}, \code{UP}, \code{B}, \code{SIM}, \code{N}, \code{RUF},
\code{ASYNC}, \code{S}; el \emph{scraper} conserva \code{E/F}) y pyright \emph{standard}.

\subsection{Capas frente a puertos y adaptadores}

Un CRUD usa capas; un servicio de integración, puertos \cite{sw:hexagonal}.
\code{core\_api} sigue \code{api} $\rightarrow$ \code{services} $\rightarrow$ \code{domain}, con
\code{infrastructure/dynamo/} como único adaptador y casos de uso sin FastAPI ni boto3.
\code{ai\_api} define en \code{domain/ports.py} protocolos (\code{LLMProvider}, \code{Embedder},
\code{Retriever}, \code{TraceLog}, \dots) que implementa \code{infrastructure/}; cambiar de
proveedor LLM toca solo esa capa y \code{api/deps.py}.

\code{Trip} es el agregado raíz (\adr{0005}): embebe días (con actividades y comidas),
alojamientos y transportes, y se lee y guarda entero. Los hijos cuelgan de
\code{/trips/\{trip\_id\}/\dots}, autorizados una vez, y \code{child\_router()} genera 25 de las
49 rutas; así se cerró una brecha real (se podían leer hijos ajenos) y
desaparecieron unas 600 líneas copiadas. \adr{0019} fijó un viaje por ciudad y una fase
(\code{upcoming}, \code{ongoing}, \code{past}) derivada de las fechas, nunca almacenada;
\code{ensure\_editable()} rechaza con 409 \code{TRIP\_LOCKED} las escrituras en viajes en curso o
pasados, y tras cada \code{PATCH} se reejecutan las invariantes de la entidad. Perfil, viaje y conversación llevan \code{version}: las escrituras son condicionales y
una carrera perdida es 409 \code{CONFLICT} (bloqueo optimista). Los servicios lanzan solo errores
de dominio, que un único manejador traduce a \code{\{message, error\_code, extras\}}; ningún
\emph{endpoint} usa \code{HTTPException}. La
identidad es un \code{Principal} (\code{subject}, \code{email}, \code{role}), ampliado en
\code{core\_api} con el UUID de la cuenta.

\subsection{Identidad y confianza entre servicios}

\code{AUTH\_MODE} elige el emisor: en \code{local}, \code{core\_api} verifica el token de Google
con \code{tokeninfo} y emite un JWT HS256 \cite{sw:rfc7519} propio; en \code{cognito}
(desplegado), un \emph{user pool} con Google como IdP emite RS256 verificado \emph{offline} contra
el JWKS \cite{sw:rfc7517}. \adr{0009} cambió el emisor de \adr{0002} (JWT con secreto compartido)
pero mantuvo la regla: \code{ai\_api} reenvía a \code{core\_api} el
\emph{bearer} del usuario, con sus mismos permisos y sin secreto entre servicios. \code{core\_api}
comprueba la cuenta en cada petición y corta al instante a un usuario desactivado; \code{ai\_api}
no, y ese usuario sigue usando la IA hasta que expira el token (60~min), riesgo aceptado en
\adr{0002}.

\subsection{Tabla única en DynamoDB}

\adr{0023} retiró el RDS PostgreSQL de \adr{0009}: costaba unos 15~€/mes, la
mitad del presupuesto de 30~€/mes; era la única razón de la VPC; y los accesos eran de
clave-valor, sin \code{join}, \code{group\_by} ni \code{ilike}, y el viaje siempre entero. La tabla
única \cite{sw:singletable} (\cref{tab:04-dynamo}) resuelve cada patrón con una \code{Query}
y cuesta menos de 1~€/mes estimado. Costes reconocidos en el ADR: la integridad pasa a la
aplicación (ítem \code{EMAIL\#} en la misma transacción, borrados en cascada a mano), no hay
migraciones (un campo ausente toma el valor por defecto de la \emph{dataclass}), las preguntas
\emph{ad hoc} exigen código o exportación, y ``otro usuario'' pasa de 403 a 404 porque la clave
incluye al propietario. \code{ai\_api} escribe en \code{<prefix>-interactions} (\adr{0024}):
cada turno es una traza con resumen (\code{DAY\#}), contexto y pasos (\code{TURN\#}), dos GSI (por
usuario y por sesión del planificador) y TTL de 90 días (\cref{sec:05-arquitectura-ia}). La traza
guarda el mensaje del usuario, el historial (hasta 8~KB) y la respuesta durante esos 90 días,
visibles para el grupo \code{admin}: un dato personal en texto libre (\cref{sec:10-analisis-critico}).

\begin{table}[tb]
\centering\footnotesize
\caption{Ítems de la tabla \code{<prefix>-core} (\adr{0023}).}
\label{tab:04-dynamo}
\begin{tabularx}{\columnwidth}{@{}lllY@{}}
\toprule
Ítem & \code{PK} & \code{SK} & Índice / nota \\
\midrule
Cuenta & \code{USER\#<id>} & \code{PROFILE} & GSI1 \code{USERS}/email \\
Email & \code{EMAIL\#<email>} & \code{EMAIL} & unicidad, misma transacción \\
Viaje & \code{USER\#<id>} & \code{TRIP\#<id>} & agregado entero; GSI2 \code{TRIPS} (admin) \\
Hilo & \code{USER\#<id>} & \code{THREAD\#<id>} & sin sus mensajes \\
Mensaje & \code{THREAD\#<id>} & \code{MSG\#<ts>\#<id>} & \emph{append-only}, orden por µs \\
\midrule
\multicolumn{4}{@{}p{\dimexpr\columnwidth}@{}}{Patrones: perfil (\code{PROFILE}); viajes o
hilos (\code{begins\_with(SK)}); hilo en orden (\code{MSG\#}); administración por GSI1/GSI2.} \\
\bottomrule
\end{tabularx}
\end{table}

\subsection{Contratos tipados}

Pydantic \cite{sw:pydantic} es la fuente de verdad: \code{just contracts} exporta
\code{docs/api/*.openapi.json} y regenera los tipos TypeScript; CI falla ante cualquier
diferencia. El planificador emite eventos SSE tipados (\adr{0015}): una unión discriminada en
\code{type} (\code{text}, \code{brief}, \code{options}, \code{itinerary\_patch}, \code{error},
\code{done}) cuyas operaciones de itinerario son otra unión en \code{op}. Ningún campo es opcional
en el cable (lo desconocido viaja como \code{null}), así que el TypeScript no tiene \code{?}. Sin
\code{response\_model} en el \emph{streaming}, \code{ai\_api/openapi.py} inyecta los modelos. El orquestador es sin estado: cada petición lleva brief, itinerario y
transcripción, y los ids \code{slot:<day>:<part>} se interpretan sin sesión.

\begin{table}[tb]
\centering\footnotesize
\caption{Tamaño y pruebas por paquete: líneas contadas en el commit \code{e7916c8}; tests ejecutados con \code{just test} el 23-09-2026, todos en verde.}
\label{tab:04-tamano}
\begin{tabularx}{\columnwidth}{@{}Yrrrr@{}}
\toprule
Paquete & LOC prod. & LOC tests & Tests & Rutas \\
\midrule
\code{travel\_common} & 829 & 414 & 37 & --- \\
\code{core\_api} & 3\,709 & 3\,111 & 199 & 49 \\
\code{ai\_api} & 10\,691 & 9\,119 & 597 & 10 \\
\midrule
Total & 15\,229 & 12\,644 & 833 & 59 \\
\bottomrule
\end{tabularx}
\end{table}

Las pruebas (\cref{tab:04-tamano}; los casos ejecutados superan a las funciones por la
parametrización) no necesitan infraestructura (moto en proceso, proveedores falsos, emisor
Cognito de prueba); los PR que tocan \emph{prompts} exigen además una prueba manual con el modelo real.

\subsection{Valoración}

Hay deuda. El doble modo de autenticación es ya permanente: duplica pruebas y configuración, y
\code{AUTH\_MODE}, \code{SECRET\_KEY} y \code{COGNITO\_*} deben coincidir a mano en dos
\code{.env} sin test que lo compruebe. La revocación de 60~min en \code{ai\_api} y la falta de
límites de uso exponen seguridad y coste. En producción solo hay un proveedor LLM,
Bedrock, sin conmutación, aunque el puerto \code{LLMProvider} lo permitiría. La documentación está desfasada: \code{docs/architecture/code-quality-review.md} describe la arquitectura SQLAlchemy
anterior a \adr{0023} sin equivalente para DynamoDB, \adr{0013} no se marcó como sustituido en
su almacenamiento y quedan directorios vacíos (\code{db/}, \code{models/},
\code{repositories/}, \code{seed/}) de la capa anterior. \Cref{sec:10-analisis-critico}
prioriza las propuestas.

\section{Arquitectura de la IA}\label{sec:05-arquitectura-ia}

El sistema es un RAG~\cite{ai:lewis2020rag} con una restricción que lo condiciona todo:
\emph{toda tarjeta procede de un documento recuperado del corpus} y los ids no recuperados se
descartan. Lugares, fotos, horarios y precios salen de un corpus por ciudad construido fuera de
línea desde fuentes abiertas y comprometido en git; el LLM interpreta la petición, elige entre ids
recuperados y redacta frases cortas. Esa prosa libre (respuestas de chat, el \code{why} de 140
caracteres) sí puede mencionar lugares y no está verificada. El turno
(\cref{fig:05-pipeline}) corre en la Lambda de \code{ai\_api} (512~MB, 900~s, streaming;
\cref{sec:06-nube}).

\begin{figure*}[tb]\centering
\begin{tikzpicture}[
  n/.style={draw, rounded corners, fill=kyrianlight, font=\footnotesize, align=center,
            minimum height=8mm, inner sep=2pt},
  a/.style={-{Stealth}, thick, kyrian},
  lab/.style={font=\footnotesize\itshape, text=kyrian},
  node distance=2.6mm]
% Fila superior: fuera de línea
\node[n, text width=3.15cm] (src) {\textbf{Fuentes abiertas}\\ Wikivoyage, Wikipedia\\ OSM, Wikidata, Meteo\\ tours, fotos};
\node[n, text width=2.9cm, right=of src] (cc) {\textbf{\code{city\_corpus}}\\ TOML de ciudad;\\ fuentes, OSM, fotos,\\ distritos, tours, clima};
\node[n, text width=1.7cm, right=of cc] (gate) {\textbf{Gate de}\\ \textbf{readiness}\\ + manifest\\ \code{cities.json}};
\node[n, text width=1.9cm, right=of gate] (jsonl) {\textbf{JSONL}\\ comprometido\\ en git\\ (25\,516 docs)};
\node[n, text width=1.55cm, right=of jsonl] (idx) {\code{just index}\\ \code{<ciudad>}};
\node[n, text width=1.6cm, right=of idx] (titan) {\textbf{Titan V2}\\ 1024 dim,\\ normalizado};
\node[n, text width=1.85cm, right=of titan, fill=accent!15] (s3) {\textbf{S3 Vectors}\\ índice \code{city-kb}\\ coseno};
\foreach \x/\y in {src/cc, cc/gate, gate/jsonl, jsonl/idx, idx/titan, titan/s3} \draw[a] (\x) -- (\y);
% Fila inferior: en línea
\node[n, text width=1.6cm, below=10mm of src.south west, anchor=north west] (req) {Petición\\ \code{PlannerTurn}\\ (estado en\\ el cliente)};
\node[n, text width=1.5cm, right=of req] (brief) {Brief\\ (LLM JSON)};
\node[n, text width=1.8cm, right=of brief, fill=accent!15] (ret) {Recuperación\\ con filtros\\ (escalera)};
\node[n, text width=1.9cm, right=of ret] (rank) {Ranking por el\\ LLM entre ids\\ recuperados};
\node[n, text width=1.6cm, right=of rank] (cards) {Tarjetas\\ hidratadas\\ del corpus};
\node[n, text width=2.1cm, right=of cards] (patch) {\code{itinerary\_}\\ \code{patch} validado\\ (reglas fijas)};
\node[n, text width=1.4cm, right=of patch] (sse) {SSE v2\\ al\\ navegador};
\node[n, text width=1.7cm, right=of sse] (trace) {Traza en\\ DynamoDB\\ (\adr{0024})};
\foreach \x/\y in {req/brief, brief/ret, ret/rank, rank/cards, cards/patch, patch/sse} \draw[a] (\x) -- (\y);
\draw[a, dashed] (patch.south) -- ++(0,-2.5mm) -| (trace.south);
\coordinate (mid) at ($(src.south)!0.5!(req.north)$);
\draw[a, dashed, accent] (s3.south) -- (s3.south |- mid) -| node[lab, pos=0.3, below] {\code{QueryVectors} / \code{GetVectors} (solo lectura)} (ret.north);
\node[lab, anchor=south west] at ($(src.north west)+(0,0.5mm)$) {Fuera de línea: portátil del operador (SSO)};
\node[lab, anchor=north west] at ($(req.south west)+(0,-4mm)$) {En línea: Lambda \code{ai\_api} (Haiku~4.5 en Bedrock, \code{eu-west-1})};
\end{tikzpicture}
\caption{Pipeline de IA. Arriba, construcción reproducible: el JSONL comprometido es la fuente de verdad y el índice se regenera con un comando. Abajo, un turno del planner.}
\label{fig:05-pipeline}
\end{figure*}

\subsection{El corpus por ciudad}

\paragraph{Fuentes y licencias.} Solo entran licencias compatibles (no Google Places,
TripAdvisor ni Booking): Wikivoyage en/es y Wikipedia, CC~BY-SA~4.0~\cite{ai:wikivoyagelicense};
OpenStreetMap (ODbL~\cite{ai:osmcopyright}) para hoteles, restaurantes, bares y distritos;
Wikidata (CC0) y Commons (sin ficheros NC, ND ni de uso legítimo) para enriquecer, y
Open-Meteo (CC~BY~4.0) para doce normales. Si \code{source}, \code{source\_url} y
\code{license} no casan, la construcción falla. Única excepción: fotos de la web del propio
hotel, con el dominio como crédito (\adr{0022}). Cada línea de \code{documents.jsonl} lleva \code{doc\_id} estable, distrito, \code{category},
\code{kind} (\code{listing}/\code{prose}), texto, coordenadas, horario, \code{price\_tier} y atribución. Los \emph{listings} son atómicos; la prosa se trocea por
encabezado en $\leq$500 tokens con un párrafo de solape.

\begin{table}[tb]\centering\footnotesize
\caption{Corpus actual (conteo propio de \code{documents.jsonl}; fecha de construcción, 2026). Hoteles: \code{sleep} con foto / total tras el filtrado de \adr{0022}; \cref{sec:02-producto} da los localizados antes de filtrar.}
\label{tab:05-corpus}
\begin{tabularx}{\columnwidth}{lrlrY}
\toprule
Ciudad & Docs & en / es & Hoteles & Fecha \\
\midrule
Berlín   & 12\,240 & 12\,187 / 53 & 341/378 & 22-09 \\
Budapest & 6\,105  & 5\,993 / 112 & 212/229 & 22-09 \\
Madrid   & 5\,574  & 5\,501 / 73  & 267/272 & 22-09 \\
Bolonia  & 1\,597  & 1\,532 / 59  & 3/118   & 19-09 \\
\midrule
Total    & 25\,516 & 98,8\,\% en  &         & \\
\bottomrule
\end{tabularx}
\end{table}

\paragraph{Ciudades, gate y manifest.} Una ciudad es un \code{cities/<slug>.toml} (borrador de
\code{just corpus-discover}) más su corpus (\cref{tab:05-corpus}); el planner no conoce ninguna
por su nombre. El \emph{gate} de \emph{readiness} define un corpus terminado: $\geq$150 lugares
de visita y $\geq$100 de comida localizados, $\geq$10 hoteles con foto, $\geq$5 distritos con
guía de barrio, 50\,\% de monumentos con foto, 12 normales y $\geq$3 tours. Los tours
son la única entrada manual (\code{curated/<slug>/tours.toml}), redactada desde la web del
operador, nunca de agregadores. El manifest \code{cities.json}
(\adr{0017}) publica alias, centro, introducción y portada: añadir una ciudad no toca Terraform. Los hoteles sin foto, cerca de la mitad (sobre todo por webs caídas), salen del
corpus (\adr{0022}). Bolonia (3/118) se construyó antes: hoy no superaría el gate de 10 hoteles
con foto y sigue publicada (los filtra el planner), riesgo aceptado y pendiente de reconstrucción. Cada tarjeta toma la imagen del
corpus, de Commons cercano, la \code{og:image} del local o un marcador (\adr{0021}). El tope de
2\,000 documentos por ciudad (TRA-210) \textbf{no está implementado}: Berlín lo supera 6,1 veces,
con un 53\,\% de restaurantes de OSM sin descripción.

\subsection{Embeddings y almacén vectorial}

Documentos y consultas se embeben con Titan Text Embeddings V2~\cite{ai:titanv2}
(1024 dimensiones, normalizado, sin asimetría consulta/documento), multilingüe: una pregunta en
español encuentra prosa inglesa. El almacén es Amazon S3 Vectors~\cite{ai:s3vectors}
(\adr{0014}): un índice \code{city-kb} (\code{float32}, 1024, coseno) creado por Terraform y
compartido, con \code{city} como metadato filtrable y clave \code{uuid5(doc\_id)} porque 1\,357
ids llevan acentos. Los metadatos filtrables (\code{category}, \code{district}, \code{price\_tier},
\code{lat}, \code{lon}\dots) ocupan como mucho 185~B de 2~KB. Se acepta búsqueda
\textbf{solo densa} y \textbf{sin filtro geográfico} (un radio son cuatro comparaciones sobre
\code{lat}/\code{lon}); la fusión BM25~\cite{ai:robertson2009bm25} +
RRF~\cite{ai:cormack2009rrf} ($k=60$) solo existe en el banco del spike.

\subsection{El spike TRA-151: Qdrant en Lambda frente a S3 Vectors}

\adr{0014} eligió S3 Vectors por simplicidad operativa, antes de medir; el spike
(\code{docs/architecture/vector-store-spike.md}, 17-09-2026) fue confirmatorio. Comparó las dos opciones viables tras
descartar \code{pgvector} en RDS (inalcanzable desde una Lambda sin VPC),
Qdrant Cloud gratuito (se suspende tras una semana
inactivo) y Qdrant en Fargate (16~€/mes expuesto, 30~€/mes privado: el doble, y ya el presupuesto íntegro de \adr{0009}). A es el
binario de Qdrant~\cite{ai:qdrant} tras Lambda Web Adapter, con la colección horneada en la imagen
y desempaquetada en \code{/tmp} en cada arranque en frío; B, S3 Vectors.

\paragraph{Método.} Ambos se llenaron con los mismos vectores de los 6\,082 documentos de
Budapest de entonces, con el coseno exhaustivo como referencia, 32 consultas y 52 \code{doc\_id}
esperados (4 en español sobre documentos en inglés). La latencia se midió desde una Lambda sin
VPC en \code{eu-west-1}, como \code{ai\_api}: 128 consultas en caliente por candidato y memoria y
arranques en frío forzados.

\begin{table}[tb]\centering\footnotesize
\caption{Calidad medida (32 consultas). Agr.: fracción del top-10 exhaustivo que devuelve el almacén.}
\label{tab:05-calidad}
\begin{tabularx}{\columnwidth}{Yrrrr}
\toprule
Estrategia & R@5 & R@10 & MRR & Agr. \\
\midrule
Coseno exhaustivo & 0,688 & 0,812 & 0,591 & --- \\
B: S3 Vectors & \textbf{0,688} & \textbf{0,812} & 0,591 & \textbf{0,969} \\
A: Qdrant en Lambda & 0,672 & 0,797 & 0,591 & 0,963 \\
Solo BM25 & 0,578 & 0,646 & 0,505 & 0,353 \\
Exhaustivo + BM25 (RRF) & 0,693 & 0,760 & \textbf{0,651} & 0,644 \\
\bottomrule
\end{tabularx}
\end{table}

\begin{table}[tb]\centering\footnotesize
\caption{Latencia medida (ms). \emph{search}: el almacén; \emph{total}: con Titan, de p50 72--75~ms en todos los grupos.}
\label{tab:05-latencia}
\setlength{\tabcolsep}{3.4pt}
\begin{tabularx}{\columnwidth}{Yrrrrrr}
\toprule
 & & \multicolumn{3}{c}{Caliente} & \multicolumn{2}{c}{Frío} \\
\cmidrule(lr){3-5}\cmidrule(lr){6-7}
Cand. & MB & search p50 & p95 & total p50 & search & total \\
\midrule
A & 2048 & 21,5 & 27,1 & 93,5 & 1\,959 & 2\,054 \\
A & 1024 & 22,5 & 29,3 & 97,7 & 1\,717 & 1\,824 \\
B & 2048 & 64,2 & 69,1 & 136,9 & 91 & 189 \\
B & 1024 & 66,5 & 152,3 & 140,1 & 288 & 387 \\
\bottomrule
\end{tabularx}
\end{table}

\begin{table}[tb]\centering\footnotesize
\caption{Coste mensual estimado por volumen (precios de lista de \code{us-east-1}, no de \code{eu-west-1}; Titan añade a ambos $\sim$0,002~USD por 100\,000 consultas).}
\label{tab:05-coste}
\begin{tabularx}{\columnwidth}{Yrr}
\toprule
Consultas/mes & A: Qdrant Lambda & B: S3 Vectors \\
\midrule
1\,000   & $\sim$0,07~USD & $\sim$0,005~USD \\
10\,000  & $\sim$0,08~USD & $\sim$0,03~USD \\
100\,000 & $\sim$0,09~USD & $\sim$0,25~USD \\
\bottomrule
\end{tabularx}
\end{table}

\paragraph{Resultados.} En calidad (\cref{tab:05-calidad}) no hay diferencia apreciable: un
documento en una consulta, con 32 consultas de una sola ciudad (4 en español) y sin intervalos de
confianza. La fusión híbrida sube el MRR de 0,591 a 0,651 y baja R@10 de 0,812 a 0,760. En caliente (\cref{tab:05-latencia}) Qdrant es tres veces más rápido, pero el
embedding domina y todo es despreciable frente a un LLM; en frío, el caso habitual con tráfico de
demostración, el orden se invierte. El coste (\cref{tab:05-coste}) es ruido frente a los 30~€/mes del presupuesto.

\paragraph{Decisión.} Se mantuvo S3 Vectors porque \emph{deciden las operaciones} y el arranque en frío: B es un índice
y un comando (6\,082 vectores en 14~s); A, una imagen de 125~MB y una colección que se reconstruye
con cada cambio del corpus y que tres veces se publicó vacía informando de éxito. Las tarjetas, que son el producto, se sirven con búsqueda densa sola: no
activar BM25 + RRF es una deuda asumida sobre lo medido, a cambio de no mantener un índice BM25 en
proceso por ciudad y de no perder R@10.

\subsection{Modelos preentrenados: por qué no hay entrenamiento ni ajuste fino}\label{sec:05-modelos}

No se entrena ni se ajusta ningún modelo: se consumen modelos fundacionales preentrenados por
API. En producción, Claude Haiku~4.5 de Anthropic en Amazon Bedrock (perfil europeo
\code{eu.anthropic.claude-haiku-4-5-20251001-v1:0}, API Converse~\cite{ai:converse}, 4\,096
tokens de salida) hace todo el trabajo generativo, y Titan Text Embeddings V2
(\code{amazon.titan-embed-text-v2:0}), los embeddings. En local, el mismo puerto habla con
\code{nvidia/nemotron-3-super-120b-a12b} por la API alojada de NVIDIA, compatible con OpenAI.
\code{eu.amazon.nova-lite-v1:0} está configurado para títulos, pero ningún código lo usa.

Se prefirió RAG con ingeniería de \emph{prompts} al ajuste fino porque el conocimiento que
importa es factual, local y cambiante (hoteles, horarios, fotos, tours): se actualiza
reconstruyendo un JSONL y reindexando por céntimos, y solo la recuperación permite exigir que
cada tarjeta cite su documento y su licencia (CC~BY-SA y ODbL piden atribución, que unos pesos
diluirían). Tampoco hay datos etiquetados de itinerarios buenos y malos ni usuarios que los
generen, y ajustar exigiría cómputo, un banco de evaluación que aún no existe
(\cref{sec:09b-evaluacion}) y repetirlo con cada modelo base: desproporcionado para cuatro personas
y 30~€/mes.

Lo que se ajusta en la práctica son los \emph{prompts} (plantillas en \code{prompts.py}, con su
\emph{hash} en la traza), la escalera de filtros y los 8 candidatos por franja
(\code{PLANNER\_CANDIDATES}), la validación determinista y la temperatura: un único
\code{CHAT\_TEMPERATURE} de 0,7, injustificado para las salidas JSON (brief, \code{classify},
ranking), donde procede 0--0,3. El JSON se pide por \emph{prompt} (esquema Pydantic en el
\code{system}), sin \emph{tool use}; si no valida hay \emph{una} llamada de reparación y, si
vuelve a fallar, el turno se degrada (mejores candidatos, títulos ``Día~N'') sin romperse. Un
ajuste fino, o un modelo pequeño propio, solo se plantearía para tareas acotadas y verificables
(extraer el brief, clasificar el mensaje) cuando las trazas (\adr{0024}) reúnan miles de turnos
reales etiquetables y el candidato supere a Haiku en el banco de evaluación con menos latencia y
coste.

\subsection{Orquestador PlanTrip}

\begin{figure}[tb]\centering
\begin{tikzpicture}[
  n/.style={draw, rounded corners, fill=kyrianlight, font=\footnotesize, align=center,
            minimum height=7mm, text width=3.4cm, inner sep=2pt},
  a/.style={-{Stealth}, thick, kyrian},
  t/.style={font=\scriptsize, text=accent},
  node distance=3.5mm]
\node[n] (b) {1. Brief: \code{extract\_brief}};
\node[n, below=of b] (nb) {3. Barrios: \code{rank\_neighbourhoods}\\ 3 tarjetas \code{nb}};
\node[n, below=of nb] (h) {4. Hoteles: \code{pick\_hotels}\\ grupo \code{hotels:<distrito>}};
\node[n, below=of h, fill=accent!15] (d) {5. Borrador: \code{skeleton} +\\ N $\times$ \code{day\_picks} + validación};
\node[n, below=of d] (c) {6. \code{classify}: opciones,\\ cambiar estancia o chat};
\node[n, right=11mm of b, text width=1.8cm] (q) {2. \code{ask\_missing}\\ (stream)};
\draw[a] (b) -- node[t, right] {brief completo} (nb);
\draw[a] (nb) -- node[t, right] {\code{select} barrio} (h);
\draw[a] (h) -- node[t, right] {\code{select} hotel} (d);
\draw[a] (d) -- node[t, right] {mensaje} (c);
\draw[{Stealth}-{Stealth}, thick, kyrian] (b) -- node[t, above] {falta} (q);
\draw[a] (c.west) -- ++(-3mm,0) |- node[t, pos=0.25, left] {\code{slot:d:p}} (d.west);
\end{tikzpicture}
\caption{Fases de PlanTrip. Sin estado en servidor: la fase se deduce del brief, del \emph{snapshot} y del \code{group\_id} de la acción.}
\label{fig:05-plantrip}
\end{figure}

\paragraph{Orquestador.} \code{PlanTrip} (\adr{0015}, \cref{fig:05-plantrip}) es \textbf{sin
estado}: cada \code{PlannerTurn} trae historial, brief e itinerario (solo ids), el
\code{group\_id} acuñado por el servidor codifica qué significa una selección y las tarjetas se
rehidratan del almacén, nunca del cliente.

\paragraph{Grounding estricto.} Cada candidato llega como una línea
\code{id | nombre | categoría | distrito | tier | foto | resumen}; el modelo solo puede usar esos
ids, con un \code{why} de menos de 140 caracteres sin precios. Los ids ajenos se
\textbf{descartan} (contados en la traza) y los huecos se rellenan con los mejores candidatos.
La tarjeta sale entera de metadatos, los precios son \emph{tiers} (€ a €€€) y \code{strip\_prices}
borra todo importe con moneda del texto del modelo. Cada
franja lanza cuatro búsquedas paralelas con una escalera de filtros (distritos del día y tier
$\rightarrow$ ciudad y tier $\rightarrow$ ciudad), deduplica por título y antepone las que tienen
foto.

\paragraph{Validación determinista.} Una función pura avisa, por día, de actividades consecutivas
a más de 3,5~km, exceso sobre el ritmo (3, 4 o 6 actividades), lugares cerrados según su horario
y precios no verificados. El tiempo sale de Open-Meteo o de la normal del
corpus; los vuelos, de un enlace sin precio.

\begin{table}[tb]\centering\footnotesize
\caption{Llamadas LLM por turno (según el código, sin reparaciones); cada búsqueda añade una a Titan y una \code{QueryVectors}.}
\label{tab:05-llamadas}
\begin{tabularx}{\columnwidth}{Yl}
\toprule
Turno & Llamadas \\
\midrule
Primer mensaje incompleto & 2 \\
Brief completo $\rightarrow$ barrios & 2 \\
Barrio elegido $\rightarrow$ hoteles & 1 \\
Hotel elegido $\rightarrow$ borrador de N días & \textbf{1 + N} ($\leq$8) \\
``Elige por mí'' con brief completo & 3 + N \\
Alternativas de una franja & 1 \\
Mensaje libre tras la estancia & 2 \\
Elegir o quitar una actividad & 0 \\
\bottomrule
\end{tabularx}
\end{table}

\paragraph{Coste.} Los días del borrador se generan en secuencia (\cref{tab:05-llamadas}) para
que el viaje tome forma en pantalla. A 1 y 5~USD por millón de tokens de entrada y salida
(Haiku~4.5) se \emph{estima} un borrador equilibrado de cuatro días en $\sim$0,02~USD
($\sim$11\,500 tokens de entrada, $\sim$2\,000 de salida) y un turno de una o dos llamadas en
0,002--0,006~USD. La traza (\adr{0024}) registra tokens, coste y latencia por turno y la consola
muestra p50/p95 (\cref{sec:09b-evaluacion}).

\paragraph{Observabilidad.} La traza de cada turno (\adr{0024}, \cref{sec:09-sre}) guarda
modelo, \emph{hash} del \emph{prompt}, tokens, coste, latencia, \emph{spans} y contadores
propios de la IA (reparaciones, ids descartados, precios eliminados); de la identidad, solo el
\emph{subject}. Los tokens de embeddings aún no se contabilizan (TRA-232).

\subsection{Valoración}

Las métricas y pruebas se recogen en la \cref{sec:09b-evaluacion}. Cada tarjeta queda anclada a un documento recuperado; la prosa libre no. Debilidades. \emph{Inyección de prompt}: el texto recuperado (Wikivoyage editable,
OSM) entra en el prompt sin delimitarse como datos y las \code{og:image} son de terceros; la
introducción española de Madrid llegó vandalizada al manifest. La
única mitigación está en la red: \code{site\_previews.py}
solo pide URL http(s) de hosts con nombre público (sin IP literales, \code{localhost} ni
\code{.local}/\code{.internal}) y revalida cada redirección contra SSRF. \emph{Capacidad}: el
límite real lo marcan las cuotas de Bedrock (no medidas), la concurrencia de Lambda y los 29~s de API
Gateway en modo \emph{buffered}; no hay pruebas de carga. Un solo proveedor y región: pasar a
NVIDIA exige un \code{apply} manual y el ciclo de vida del modelo lo fija el proveedor. Además, un
corpus 98,8\,\% en inglés ante usuarios hispanohablantes; la
escalera reembebe la consulta en cada escalón, sin caché; y no hay reranker ni guardrails más
allá de \code{strip\_prices} (\cref{sec:10-analisis-critico}).


\section{Arquitectura en la nube}\label{sec:06-nube}

El sistema vive en una cuenta de AWS (\code{eu-west-1}), sin VPC desde el 23 de septiembre de
2026. Todo está en Terraform (proveedor \code{hashicorp/aws \textasciitilde{}> 6.26}, estado en S3
con \code{use\_lockfile}) salvo el
\emph{bootstrap} (bucket de estado, proveedor OIDC y rol de CI, aplicado una vez con estado local) y
el WAF, creado desde la consola. \Cref{fig:06-aws} muestra la topología \emph{v3}.

\begin{figure*}[tb]
\centering
\begin{tikzpicture}[
  n/.style={draw, rounded corners, fill=kyrianlight, font=\footnotesize, align=center,
            minimum height=6mm, text width=2.6cm, inner sep=1pt},
  ext/.style={n, fill=white, dashed},
  a/.style={-{Stealth}, thick, draw=kyrian},
  d/.style={-{Stealth}, dashed, draw=gray},
  l/.style={font=\scriptsize, inner sep=1.5pt},
  grp/.style={draw=kyrian!60, rounded corners=3pt, inner sep=4pt},
  gl/.style={font=\scriptsize\bfseries\color{kyrian}, anchor=south west, inner sep=1pt},
  x=1cm, y=0.74cm]
  % nodos
  \node[n, text width=1.7cm, fill=white] (user) at (-0.3,0.4) {Usuario\\(navegador)};
  \node[n] (cognito) at (6.6,3.3) {Cognito User Pool\\RS256, grupo \code{admin}};
  \node[n] (google) at (10.3,3.3) {Google\\(único IdP)};
  \node[n] (cf) at (2.9,0.4) {CloudFront + WAF\\ACM, Route~53};
  \node[n] (s3) at (2.9,-1.2) {S3 privado (OAC)\\export estático};
  \node[n] (apigw) at (6.6,0.4) {API Gateway REST\\autorizador Cognito};
  \node[n] (core) at (10.3,1.2) {Lambda \code{core-api}\\{\scriptsize \emph{buffered}, 1769~MB, 30~s}};
  \node[n] (ai) at (10.3,-0.5) {Lambda \code{ai-api}\\{\scriptsize SSE, 512~MB, 900~s}};
  \node[n] (ddbc) at (14.3,2.25) {DynamoDB \code{core}\\PITR, protección};
  \node[n] (ddbi) at (14.3,1.1) {DynamoDB\\\code{interactions}, TTL};
  \node[n] (vec) at (14.3,-0.05) {S3 Vectors\\índice \code{city-kb}};
  \node[n] (bed) at (14.3,-1.25) {Bedrock: Haiku~4.5,\\Nova Lite, Titan V2};
  \node[ext] (ofm) at (-0.1,-2.8) {OpenFreeMap\\(teselas del mapa)};
  \node[ext] (om) at (14.3,-2.8) {Open-Meteo\\(previsión)};
  \node[n, text width=2.3cm] (gha) at (6.6,-2.8) {GitHub Actions\\OIDC, sin claves};
  \node[n, text width=2.1cm] (ecr) at (10.3,-2.8) {ECR\\(imágenes)};
  % grupos
  \begin{scope}[on background layer]
    \node[grp, fit=(cf)(s3)] (gb) {}; \node[gl] at (gb.north west) {borde};
    \node[grp, fit=(apigw)(core)(ai)] (gc) {}; \node[gl] at (gc.north west) {cómputo};
    \node[grp, fit=(ddbc)(bed)] (gd) {}; \node[gl] at (gd.north west) {datos e IA};
    \node[grp, fit=(google)(cognito)] (gi) {}; \node[gl] at (gi.north west) {identidad};
    \node[grp, dashed, fit=(ofm)] (ge1) {}; \node[gl] at (ge1.north west) {externos};
    \node[grp, dashed, fit=(om)] (ge2) {}; \node[gl] at (ge2.north west) {externos};
    \node[grp, fit=(gha)(ecr)] (gci) {}; \node[gl] at (gci.north west) {entrega};
  \end{scope}
  % flujos
  \draw[a] (user) -- node[l, above]{HTTPS} (cf);
  \draw[a] (cf) -- node[l, right]{\code{default}} (s3);
  \draw[a] (cf) -- node[l, above]{\code{/api/*}} (apigw);
  \draw[a] (apigw.east) -- (core.west);
  \draw[a] (apigw.east) -- (ai.west);
  \draw[a] (core.east) -- (ddbc.west);
  \draw[a] (ai.east) -- (ddbi.west);
  \draw[a] (ai.east) -- (vec.west);
  \draw[a] (ai.east) -- (bed.west);
  \draw[d] (ai.south east) -- (om.north west);
  \draw[d] (user.north) |- node[l, pos=0.75, above]{inicio de sesión: \code{code}+PKCE} (cognito.west);
  \draw[d] (cognito) -- (google);
  \draw[d] (apigw) -- node[l, right, pos=0.35]{valida JWT} (cognito);
  \draw[d] (user) -- node[l, right]{mapa} (ofm);
  \draw[a] (gha) -- node[l, below]{\code{crane copy}} (ecr);
  \draw[a] (ecr) -- node[l, right]{imagen por digest} (ai);
\end{tikzpicture}
\caption{AWS v3 en \code{eu-west-1}, sin VPC. API Gateway exige un ID token de Cognito en todo
método (incluidos los de salud). Ambas funciones reciben su imagen de ECR por digest (solo se dibuja
la flecha a \code{ai-api}). Discontinuo: identidad y servicios externos. Fuente: \code{infra/aws/*.tf}
y \code{docs/architecture/aws-architecture.drawio.svg}.}
\label{fig:06-aws}
\end{figure*}

\subsection{Topología v3}

\textbf{Borde.} CloudFront sirve el export desde S3 privado (versionado, solo \emph{Origin
Access Control} con \code{SourceArn}) y reenvía \code{/api/*} sin caché: un único origen, sin
CORS. Una CloudFront Function traduce
\code{/x/} a \code{/x/index.html}; \code{/api/*} va sin compresión para que el SSE llegue token a
token.

\textbf{Cómputo.} API Gateway REST regional tiene dos \emph{proxy}: \code{/api/v1/\{proxy+\}} hacia
\code{core-api} en modo \emph{buffered} (límite de 29~s, de ahí el \emph{timeout} de 30~s) y
\code{/api/v1/ai/\{proxy+\}} hacia \code{ai-api} en \code{STREAM} (hasta 15~min). Un
\code{Dockerfile} construye ambas imágenes; Lambda Web Adapter~\cite{aws:lwa} solo se activa bajo
Lambda, y el contenedor corre igual en Compose (\cref{sec:07-entorno-desarrollo}). \code{core-api} tiene 1769~MB (una vCPU);
\code{ai-api}, que casi solo espera al modelo, 512~MB. Sin \emph{provisioned concurrency}.

\textbf{Datos e IA.} DynamoDB \emph{on-demand}: \code{core}
(\cref{sec:04-arquitectura-software}) e \code{interactions} (\cref{sec:09-sre}). Los
vectores están en S3 Vectors~\cite{aws:s3vectors} (índice \code{city-kb}, 1024 dimensiones,
coseno) y los modelos en Bedrock con perfiles de inferencia EU: Claude Haiku~4.5 responde y Titan
V2 embebe consultas; Nova Lite está declarado para títulos, pero el código no lo invoca
(\cref{sec:05-arquitectura-ia}).

\textbf{Identidad.} Cognito, solo con Google y sin registro por contraseña
(\code{allow\_admin\_create\_user\_only}), emite JWT RS256 (\emph{ID} y \emph{access} de 60~min,
\emph{refresh} de 30~días) con \code{code}+PKCE y cliente público. Terraform lee el JWKS en \code{plan} y lo
inyecta como variable de entorno: validar un token no sale a la red. Sin VPC no hace falta NAT:
las Lambdas salen a Internet por la red gestionada de AWS (Open-Meteo, Commons, \code{og:image}),
y Bedrock, DynamoDB y S3 Vectors se autorizan por rol IAM.

\subsection{Evolución v1\,$\rightarrow$\,v3}

Tres iteraciones en ocho días (15--23 de septiembre; \cref{tab:06-evolucion}) con un
presupuesto real de 60~€ para dos meses: la v2 costaba \mbox{$\sim$55}~€/mes antes de resolver
la salida a Internet (NAT \mbox{$\sim$35}~€/mes o cinco \emph{interface endpoints}
\mbox{$\sim$40}~€/mes). ECS no era requisito: Lambda sustituyó a Fargate a costa de no tener tareas en segundo plano
ni estado en memoria entre peticiones.

\begin{table}[tb]
\centering\footnotesize
\caption{Evolución de la arquitectura AWS.}\label{tab:06-evolucion}
\begin{tabularx}{\columnwidth}{@{}>{\raggedright\arraybackslash}p{1.45cm}YY@{}}
\toprule
Versión & Qué cambió & Por qué \\
\midrule
v1 (\adr{0007}) & Fargate + ALB + RDS; SSO local, OIDC en CI & Sin credenciales largas legibles por agentes \\
v2 (\adr{0008}) & CloudFront origen único; API Gateway con \emph{streaming}; VPC Link + NLB; Bedrock & ALB HTTP bloqueaba Google Sign-In; CORS. NAT y NLB sin aplicar \\
v3 (\adr{0009}) & Lambda + LWA; Cognito; sin NAT; VPC solo por RDS & 30~€/mes; sin salida (\code{tokeninfo}, NVIDIA) \\
sin VPC (\adr{0023}) & DynamoDB; RDS y VPC destruidos (TRA-219) & RDS: $\sim$15~€/mes y la VPC \\
\bottomrule
\end{tabularx}
\end{table}

\subsection{Costes}

\Cref{tab:06-costes} son estimaciones de los ADR (precios de lista, tráfico de
demostración~\cite{aws:pricing}): \mbox{$\sim$5}~€/mes en recursos gestionados por Terraform
(\adr{0023}) más \mbox{$\sim$8}~\$/mes del WAF creado a mano, unos 12--13~€/mes en total. No hay
factura en esta memoria: la de septiembre debe sustituir la estimación.

\begin{table}[tb]
\centering\footnotesize
\caption{Coste mensual estimado (€/mes) antes y después de retirar RDS.}\label{tab:06-costes}
\begin{tabularx}{\columnwidth}{@{}Yccc@{}}
\toprule
Componente & \adr{0009} & \adr{0023} & Escala \\
\midrule
Lambda $\times$2 (capa gratuita) & 0 & 0 & sí \\
API Gateway (3,50~\$/M pet.) & $<$1 & $<$1 & sí \\
CloudFront, S3, Route~53 & $\sim$1 & $\sim$1 & $>$1~TB \\
Cognito (10\,000 MAU gratis) & 0 & 0 & $>$10\,000 \\
RDS \code{db.t4g.micro} + 20~GB & $\sim$15\textsuperscript{a} & 0 & no \\
DynamoDB (2 tablas) + PITR & --- & $<$1 & sí \\
ECR, CloudWatch Logs, KMS & $\sim$1 & $\sim$1 & logs \\
Bedrock (tokens, demo) & 1--3 & 1--3 & sí \\
\midrule
\textbf{Total Terraform} & \textbf{$\sim$4 / $\sim$19} & \textbf{$\sim$5} & \\
WAF (consola, fuera de Terraform) & & $\sim$8~\$ & no \\
\textbf{Total actual} & & \textbf{$\sim$12--13} & \\
\bottomrule
\multicolumn{4}{@{}p{\columnwidth}@{}}{\scriptsize \textsuperscript{a}0~€ el primer año (capa
gratuita).}
\end{tabularx}
\end{table}

\subsection{Seguridad}

Sin credenciales de larga duración (\adr{0007} prohíbe \emph{access keys} en \code{.env},
\code{tfvars}, repositorio y contenedor): en local, IAM Identity Center (\code{just aws-login}); en
CI, OIDC con un rol de 1~h que solo confía en el entorno \code{aws}. Roles mínimos: \code{core-api} sin \code{Scan}; \code{ai-api} con sus trazas, un índice de
vectores de solo lectura y Bedrock limitado a perfiles EU. CI parte
de \code{PowerUserAccess} con IAM acotado a \code{travel-ai-*}. PITR y
\emph{deletion protection} cubren \code{core} y el \emph{user pool}; el WAF, tres grupos de
reglas gestionadas. Contra SSRF, \code{site\_previews.py} solo pide \code{og:image} a hosts con nombre
público y revalida cada redirección (\cref{sec:05-arquitectura-ia}). Carencias: sin \emph{usage plans} ni \emph{throttling}, cualquier usuario
autenticado puede generar gasto ilimitado en Bedrock; el estado de Terraform
contiene dos secretos (clave NVIDIA y secreto OAuth de Google); en local se usa
\code{AdministratorAccess} (la organización no permitió un \emph{permission set} acotado); no hay
GuardDuty, Config ni Security Hub.

\subsection{Raíces del dominio y la alternativa GCP}

\adr{0010} importó al estado, con \code{prevent\_destroy}, la zona de Route~53 y el certificado ACM
(\code{us-east-1}, como exige CloudFront), antes \emph{data sources}: un \code{apply} reconstruye el borde sin exponer
recursos irreversibles a un \code{destroy}. \code{infra/gcp/} solo pasa \code{terraform validate} en CI y, según su README, no está
portado a DynamoDB, así que no despliega \code{core\_api}.

\subsection{Valoración}

Cada ADR corrigió al anterior con cifras de la cuenta real y el coste fijo estimado pasó de
\mbox{$\sim$55}~€ a \mbox{$\sim$12--13}~€/mes (\mbox{$\sim$5}~€ en Terraform más el WAF). Deuda:
WAF fuera de IaC y del presupuesto; una sola región, sin \emph{staging}; los
despliegues fijan el digest sin alias ni tráfico ponderado; arranques en frío de 1--3~s sin mitigar; y \code{interactions} no tiene PITR ni protección contra borrado: coherente para un log con TTL,
pero deja sin copia el único registro de lo que hizo la IA. \code{infra/gcp/} debe retirarse o
portarse (\cref{sec:10-analisis-critico}).

\section{Entorno de desarrollo}\label{sec:07-entorno-desarrollo}

Como desarrollan a la vez una persona y varios agentes con \emph{shell}, el \emph{devcontainer}~\cite{ci:devcontainers} (base Ubuntu de Microsoft) instala las versiones de CI: Python~3.12 con
\code{uv}, Node.js~24, \code{just}, Terraform~1.16 (\emph{build arg}), AWS CLI~v2,
\code{crane}~0.22.1, \code{gh}, \code{ripgrep}/\code{jq}, Claude Code y, opcionales, Codex,
Gemini y Copilot CLI. \code{post-create.sh} copia un perfil SSO de ejemplo y ejecuta
\code{just setup} (cuatro \code{.env}, \code{uv sync}, \code{npm install}) y Chromium de
Playwright. Compose añade DynamoDB Local y volúmenes que conservan \code{.venv}, \code{node\_modules}, la configuración de los agentes
y \code{\textasciitilde/.aws}. AWS se usa por IAM Identity Center (\code{just aws-login}),
sin claves de larga duración en disco que un agente pudiera leer (\adr{0007}). Se descartaron Docker-in-Docker, SAM y CDK; el \emph{stack} de Compose se lanza
desde el \emph{host} porque monta rutas relativas a su sistema de ficheros.

Las 41 recetas del \code{justfile} (\cref{tab:07-just}) son la única interfaz para personas,
agentes y los \emph{jobs} de \code{pr.yml}. Con \code{dynamodb-local} (moto, \code{:8002}),
tests y desarrollo no necesitan Docker. \code{stack-up} sirve \emph{export}, APIs y DynamoDB Local en un origen
(\code{:8080}) tras un nginx que reproduce CloudFront (SSE sin \emph{buffering}, \code{/api/} a \code{core\_api}, índices y \code{404.html}); lo ejecuta el \emph{job}
\code{e2e-stack} (\cref{tab:07-paridad}).

\begin{table}[tb]
\centering\footnotesize
\caption{Recetas de \code{just} agrupadas por uso (selección de las 41).}
\label{tab:07-just}
\begin{tabularx}{\columnwidth}{@{}lY@{}}
\toprule
Grupo & Recetas \\
\midrule
Desarrollo & \code{setup}, \code{dev-\{core,ai,frontend\}}, \code{dynamodb-local} \\
Calidad & \code{lint}, \code{test\{,-core,-ai,-common,-corpus,-frontend\}}, \code{test-e2e\{,-static,-stack\}}, \code{docs-check}, \code{contracts\{,-check\}} \\
Corpus & \code{corpus-discover}, \code{corpus}, \code{corpus-report}, \code{index} \\
\emph{Stack} y nube & \code{docker-up}, \code{build-stack}, \code{stack-up}, \code{aws-login}, \code{infra-\{validate,fmt-check\}}, \code{release} \\
\bottomrule
\end{tabularx}
\end{table}

\begin{table}[tb]
\centering\footnotesize
\caption{Paridad entre el entorno local y AWS.}
\label{tab:07-paridad}
\begin{tabularx}{\columnwidth}{@{}lYl@{}}
\toprule
Aspecto & Local frente a AWS & Paridad \\
\midrule
Cómputo & misma imagen; uvicorn / Lambda Web Adapter & alta \\
Origen & nginx \code{:8080} / CloudFront & alta \\
Datos & DynamoDB Local o moto / DynamoDB & alta \\
LLM & NVIDIA NIM con clave / Bedrock por rol IAM & media \\
Auth & \code{AUTH\_MODE=local} (JWT HS256 por CLI) / Cognito RS256 & media \\
Vectores & sin emulador: índice real de S3 Vectors, requiere SSO & baja \\
WAF & sin equivalente local & nula \\
WebGL & SwiftShader sin interfaz; quitar \code{WAYLAND\_DISPLAY} (TRA-180) & aviso \\
\bottomrule
\end{tabularx}
\end{table}

\paragraph{Valoración.} Un agente recién arrancado ejecuta las comprobaciones de CI sin más
instrucciones, y cada PR prueba el \emph{export} como lo sirve CloudFront. La paridad cae donde
hay un servicio gestionado sin emulador: el RAG local depende de AWS real y no funciona sin conexión, y el WAF solo se comprueba
en producción. El \emph{stack} completo no corre dentro del devcontainer, que arrastra herramientas sin
uso (el \emph{plugin} de Session Manager de la etapa Fargate).

\section{CI/CD y proceso de entrega}\label{sec:08-cicd-proceso}

Cada cambio empieza en un issue de Linear y llega a producción por dos caminos
(\cref{fig:08-flujo}): el frontend se despliega al fusionar; el backend, a mano.

\begin{figure*}[tb]
\centering
\begin{tikzpicture}[
  n/.style={draw, rounded corners, fill=kyrianlight, font=\footnotesize, align=center,
            minimum height=8mm, text width=2.6cm, inner sep=2pt},
  h/.style={n, fill=accent!15},
  a/.style={-{Stealth}, semithick},
  l/.style={font=\scriptsize, fill=white, inner sep=1pt}]
% fila 1: la ola de agentes
\node[n] (iss) at (0,0)     {Issue Linear\\\code{TRA-<n>}, aceptación};
\node[n] (bri) at (3.55,0)  {\textbf{Brief}\\sesión principal\\(modelo más capaz)};
\node[n] (imp) at (7.1,0)   {\textbf{Implementar}\\Opus, \emph{worktree}\\rama \code{<tipo>/TRA-<n>-...}};
\node[n] (rev) at (10.65,0) {\textbf{Revisar} 3$\times$Sonnet\\corrección $\cdot$ contratos\\y tests $\cdot$ UX/a11y};
\node[n] (fix) at (14.2,0)  {\textbf{Corregir}\\Opus: \emph{findings},\\\code{gh pr checks --watch}};
\draw[a] (iss) -- (bri); \draw[a] (bri) -- (imp);
\draw[a] (imp) -- node[l,above=1pt]{PR} (rev); \draw[a] (rev) -- node[l,above=1pt]{JSON} (fix);
% fila 2: CI de PR y merge
\node[n, text width=9.7cm] (pr) at (10.65,-1.75) {\code{pr.yml} (\emph{jobs} filtrados por ruta): lint (ruff, pyright, eslint) $\cdot$
  tests por paquete $\cdot$ \code{docker-build} sin publicar $\cdot$ \emph{drift} de contratos $\cdot$
  e2e del \emph{stack} $\cdot$ \code{terraform validate} $\cdot$ \code{docs-check}};
\node[h] (mer) at (0,-1.75) {\textbf{Merge} \emph{squash}\\a \code{main} (humano\\o sesión principal)};
\draw[a] (imp.south) -- node[l,right]{abre} (imp.south|-pr.north);
\draw[a] (fix.south) -- node[l,right]{\emph{push}} (fix.south|-pr.north);
\draw[a] (pr.west) -- node[l,above=1pt]{verde} (mer.east);
% fila 3: backend
\node[n] (img) at (0,-3.5)    {\code{backend-images.yml}\\build y \emph{push} a GHCR};
\node[h, text width=6.4cm] (dep) at (5.325,-3.5) {\code{deploy-backend.yml} (\code{workflow\_dispatch}): OIDC $\to$
  \code{crane copy} GHCR$\to$ECR $\to$ \emph{digest} $\to$ \code{plan}; \code{apply} solo con \code{apply=true}};
\node[n] (lam) at (10.65,-3.5) {2 Lambda\\(imagen por \emph{digest})};
% fila 4: frontend
\node[n] (fro) at (0,-5)      {\code{deploy.yml}\\\emph{export} de Next.js};
\node[n, text width=6.4cm] (s3) at (5.325,-5) {OIDC $\to$ \code{s3 sync --delete} $\to$ invalidación de CloudFront \code{/*}};
\draw[a] (mer) -- node[l,right]{\code{src/backend/**}} (img);
\draw[a] (img) -- (dep); \draw[a] (dep) -- (lam);
\draw[a] (mer.west) -- ++(-0.2,0) |- (fro.west);
\draw[a] (fro) -- node[l,above=1pt]{\code{src/frontend/**}} (s3);
\node[font=\scriptsize, align=left, anchor=west, text width=5.6cm] at (9.3,-5)
  {\textcolor{accent}{\rule{2mm}{2mm}} decisión humana explícita\\(\emph{merge} si el clasificador lo\\deniega; \code{terraform apply}).};
\end{tikzpicture}
\caption{Del issue a producción: ola de agentes (fila superior), CI de la PR, imágenes y despliegue
del backend (manual) y del frontend (automático al fusionar).}
\label{fig:08-flujo}
\end{figure*}

\subsection{Integración continua}
\code{pr.yml} corre en cada PR contra \code{main}. Un \emph{job} inicial evalúa con
\code{dorny/paths-filter} diez filtros de ruta (uno por paquete, más \code{stack}, \code{infra} y
\code{docs}) y cada \emph{job} depende del suyo: un cambio de documentación no ejecuta pytest ni
Playwright. La \cref{tab:08-pruebas} resume los niveles de prueba: \code{contracts-check}
regenera el OpenAPI y los tipos TypeScript y falla si difieren de lo comprometido, y
\code{e2e-stack} ejecuta Playwright sobre el \emph{stack} de Compose
(\cref{sec:07-entorno-desarrollo}) con una sesión generada en \code{core\_api}. CI construye además
cada imagen, sin \code{provenance} ni \code{sbom}, para comprobar el Lambda Web Adapter sin
publicarla. Terraform solo pasa \code{fmt} y \code{validate}: ninguna PR ejecuta un \code{plan}.

\begin{table}[tb]
\centering\footnotesize
\caption{Plan de pruebas. Casos ejecutados con \code{just test} el 23-09-2026, todos en verde
(\code{travel\_common} 37, \code{core\_api} 199, \code{ai\_api} 597, \code{city\_corpus} 254;
muchos tests están parametrizados). L: local.}
\label{tab:08-pruebas}
\begin{tabularx}{\columnwidth}{@{}l>{\raggedright\arraybackslash}p{2.15cm}lY@{}}
\toprule
Nivel & Herramienta & Dónde & Casos / cobertura \\
\midrule
Unitario & pytest; Vitest & L, CI & 1\,087 backend; 747 frontend \\
Integración & pytest, moto, LLM guionizado & L, CI & incluidos en los 1\,087 \\
Contrato & \code{contracts-check} & CI & \emph{drift} OpenAPI y tipos \\
E2E & Playwright (\emph{export}, Compose) & CI & 9 \emph{specs}, no ejecutados aquí \\
Modelo real & \code{planner-smoke} (NVIDIA) & L & manual; no contra Bedrock \\
Seguridad & reglas \code{S} de ruff & L, CI & estático; sin DAST ni \emph{pentest} \\
Carga, a11y & --- & --- & \textbf{no existen} (a11y: solo revisor UX) \\
\bottomrule
\end{tabularx}
\end{table}

\subsection{Despliegue}
Al fusionar, \code{backend-images.yml} publica \code{core-api} y \code{ai-api} en GHCR,
etiquetadas con el SHA completo. \code{deploy-backend.yml} solo se lanza a mano, en el entorno \code{aws}
de GitHub, que guarda los dos únicos secretos reales (clave de NVIDIA y secreto del cliente de
Google); entra en AWS por OIDC~\cite{ci:gh-oidc}, sin claves, copia las imágenes a ECR con
\code{crane copy}~\cite{ci:crane} y pasa el \emph{digest} a Terraform. Se descartó
\code{docker buildx imagetools create} porque envuelve la imagen en un índice OCI que Lambda
rechaza, y \code{crane} no necesita \emph{daemon} (TRA-133); el \emph{digest}, a diferencia de una
etiqueta mutable, obliga a Lambda a redesplegar. El \code{apply} es manual (\code{apply=false} por
defecto): solo hay un entorno, sin \emph{staging}, y un secreto ausente llega a Terraform como
cadena vacía y borra la configuración viva (lo advierte el propio \emph{workflow}), así que el
\emph{runbook} exige revisar antes un \code{plan} que solo cambie los dos \code{image\_uri}. El
frontend (\code{deploy.yml}) no tiene esa puerta: cada \emph{push} a \code{main} que toca
\code{src/frontend/**} se sincroniza con S3 e invalida CloudFront.

Dependabot abre cada semana PRs agrupadas y retiene, con el motivo escrito, tres versiones
mayores (TRA-195).

\paragraph{Entregables: imágenes Docker y scripts de despliegue.} Un único
\code{src/backend/Dockerfile} multietapa produce las dos imágenes con el \emph{build arg}
\code{SERVICE}: \code{uv sync --frozen} instala solo ese servicio y \code{travel\_common}, y la etapa
final (\code{python:3.12-slim}, usuario sin privilegios, \code{HEALTHCHECK}) lleva el Lambda Web
Adapter como extensión, activa solo bajo Lambda; fuera, \code{docker/entrypoint.sh} arranca uvicorn.
La misma imagen corre así en Compose y en Lambda. \code{docker-compose.yml} define \code{proxy}
(nginx, \code{docker/nginx.conf}), \code{core\_api}, \code{ai\_api} y DynamoDB Local. El frontend no
tiene imagen: es un \emph{export} estático, sin proceso servidor, que sirven nginx o S3+CloudFront.
Se arranca con \code{just setup}, \code{docker-up} y \code{stack-up}, y se despliega con los
\emph{workflows} y \code{terraform apply} (\emph{runbooks} \code{deploy.md} y \code{release.md}).

\subsection{Entrega con agentes}
Según \code{AGENTS.md}, se entrega con agentes salvo documentación y correcciones de una línea,
en cuatro roles. (1)~La sesión principal, con el
modelo más capaz, audita el código, decide y escribe un \emph{brief} de secciones fijas
(decisiones que no se reabren, tests, documentación, reglas); resultado y lista de aceptación se
copian al issue, el registro permanente. (2)~\code{/deliver-issue} lanza
\code{.claude/workflows/kyrian-wave.js}: un agente Opus~5.5 implementa en su \emph{git worktree},
ejecuta los \emph{checks} y abre la PR con \code{Closes TRA-<n>}. (3)~Tres revisores Sonnet~5 en
paralelo, uno por lente (corrección y estado; contratos, reglas y tests reales; UX, textos y
accesibilidad a 390~px en los dos idiomas), devuelven \emph{findings} JSON con fichero, línea y
severidad (\code{blocker}, \code{should}, \code{nit}). (4)~Otro Opus aplica los que confirma,
justifica los que omite y vigila CI con hasta dos rondas de corrección; tiene prohibido fusionar.
Fusiona la sesión principal o, si el clasificador de permisos lo deniega, el humano, que además
ejecuta siempre el \code{terraform apply}.

Lo sostienen un \code{AGENTS.md} jerárquico, fuente única para
cualquier herramienta (Claude Code~\cite{ci:claudecode}, Codex, Gemini); diez comandos
\emph{slash}; un \emph{hook}
\code{PreToolUse} que deniega ramas fuera de \code{<tipo>/TRA-<n>-<slug>}, con los prefijos de
Conventional Commits~\cite{ci:conventional}; los MCP de Linear y Playwright; 23 ADR al estilo de
Nygard~\cite{ci:nygard}, nunca borradas, solo sustituidas; 7 \emph{runbooks}, y reglas no negociables.

\begin{table}[tb]
\centering\footnotesize
\caption{Hitos del proyecto (\code{git log}, 2026).}
\label{tab:08-hitos}
\begin{tabularx}{\columnwidth}{@{}lY@{}}
\toprule
Fecha & Hito \\
\midrule
10/18-03 & Primer \emph{commit} (landing); primer issue citado (TRA-28) \\
jun--ago & 6 \emph{commits} en tres meses \\
07-09 & Separación en \code{core\_api}/\code{ai\_api} (\adr{0001}) \\
15/16-09 & AWS v3 con Lambda y Cognito (\adr{0009}); \code{crane} (TRA-133) \\
18-09 & \emph{Spike} de vectores (medido 17-09, PR \#118); planificador \code{/plan/} \\
19-09 & Bolonia, primera ciudad con \code{/add-city} \\
20-09 & Mapa por día; cambio de marca a Kyrian World \\
22/23-09 & DynamoDB sustituye a RDS; VPC retirada (\adr{0023}) \\
23-09 & Trazas de turno y consola \code{/admin} (\adr{0024}) \\
\bottomrule
\end{tabularx}
\end{table}

\paragraph{Cifras.} Contadas con \code{git} y \code{gh} el 23-09-2026: 324 \emph{commits}, 181 PR
fusionadas y 107 claves \code{TRA-<n>} distintas. Septiembre concentra 145 PR (80\,\%) y 187
\emph{commits} (58\,\%), frente a 36 PR entre marzo y agosto. El salto coincide con la adopción
de la entrega con agentes, pero también con la construcción del planificador
(\cref{tab:08-hitos}); los datos no permiten atribuirlo a una causa. Miden volumen, no valor (faltan
\emph{reverts}, proporción \code{fix}/\code{feat} y CI fallida).

\paragraph{Valoración crítica.} A favor: toda rama, PR y \emph{commit} se traza hasta un issue y,
si procede, una ADR; el \emph{drift} de contratos se detecta en CI y no en producción; y las tres
lentes evitan que un agente valide su propio trabajo. En contra, con igual peso: \emph{bus factor}
1 (145 de las 181 PR son de la misma persona, la única que fusiona); revisión sin humano, con tres
revisores del mismo modelo y sesgos correlacionados, sin segunda revisión tras la corrección (solo
CI) ni QA manual sistemático; un coste en \emph{tokens} por ola sin medir; deriva documental; y
ningún \emph{smoke test} contra AWS tras desplegar, ni \code{plan} en la PR ni \emph{staging}: el
primer contacto con producción es el despliegue mismo.

\section{Fiabilidad, observabilidad y costes}\label{sec:09-sre}

Se siguen AWS Well-Architected~\cite{aws:wellarchitected} y la terminología de SLO de
Google~\cite{aws:srebook}; las ausencias se comprobaron en \code{infra/aws/*.tf} y en CI.

\subsection{Observabilidad: lo que hay y lo que falta}

\textbf{Existe} el sistema de monitorización y registro con seguimiento del rendimiento y
trazabilidad del flujo de datos que pide el guion, en tres capas. (1)~\emph{Registro}: cada función
escribe en su grupo de CloudWatch Logs (14~días, sin exportación) texto con pares clave=valor, no
JSON (\code{travel\_common.http.logging}): la línea de acceso de uvicorn por petición, modelo y
tokens de Bedrock por respuesta, reintentos y fallos de dependencias, y una línea de auditoría
\code{admin\_read} (sujeto, ruta, objetivo) por cada lectura administrativa. (2)~\emph{Trazabilidad}:
la traza de turno (\adr{0024}) guarda por petición en la tabla \code{interactions} (TTL de
90~días) un resumen y sus pasos (recuperaciones con su top-$k$ y distancias, llamadas al modelo,
herramientas, eventos SSE) con tokens, coste según una tabla de precios versionada y latencia por
paso, en vocabulario de OpenTelemetry GenAI, enlazada a usuario, sesión y viaje: se recorre
petición $\to$ recuperación $\to$ modelo $\to$ eventos $\to$ respuesta. Es \emph{best effort} y se
escribe antes del \code{[DONE]}, porque Lambda puede congelar el proceso al cerrarse el
\emph{stream}. (3)~\emph{Rendimiento}: la consola \code{/admin/} agrega por día turnos, errores,
tokens, coste, latencia p50/p95, primer evento, usuarios y sesiones, e inspecciona cada turno.
La traza guarda texto libre (mensaje, historial y respuesta) y no se borra al suprimir la cuenta:
un problema de privacidad (RGPD) pendiente (\cref{sec:10-analisis-critico}). Los \emph{endpoints}
de salud sirven de \emph{readiness} a Lambda Web Adapter, pero quedan tras el autorizador de
Cognito.

\textbf{No existe.} Ni alarmas de CloudWatch, SNS, AWS Budgets, dashboards, X-Ray ni colector
OTel, ni comprobación externa de disponibilidad (imposible sin credenciales), ni SLO ni
presupuesto de errores, ni \emph{smoke test} contra producción tras desplegar: \code{e2e-stack}
prueba Compose en CI y \code{deploy.md} solo pide comparar el \emph{digest}. Una caída o un
pico de gasto pasan inadvertidos hasta que alguien los ve.

\subsection{Dependencias y alta disponibilidad}

\begin{table}[tb]
\centering\footnotesize
\caption{Dependencias externas y su efecto en la disponibilidad.}\label{tab:09-dependencias}
\begin{tabularx}{\columnwidth}{@{}>{\raggedright\arraybackslash}p{1.95cm}YY@{}}
\toprule
Dependencia & Modo de fallo & Mitigación actual \\
\midrule
Bedrock (Haiku, Titan) & Caen chat, planificador y recuperación & 3 intentos (1 y 2~s) antes del primer \emph{token}; NVIDIA manual (\code{apply}) \\
DynamoDB \code{core} & Cuentas, viajes y chats inaccesibles & Multi-AZ gestionado; PITR \\
DynamoDB \code{interactions} & Se pierden trazas & \emph{Best effort}; nunca rompe la respuesta \\
S3 Vectors & Sin candidatos para el planificador & Ninguna; \code{just index} lo reconstruye \\
Cognito & Sin sesión; la API lo rechaza todo & Ninguna \\
Google & Sin nuevos inicios de sesión & Ninguna (único IdP) \\
Open-Meteo & Sin previsión & \emph{Timeout} de 5~s; normales climáticas del corpus \\
OpenFreeMap & Mapa sin teselas & Ninguna; el planificador sigue \\
GitHub / GHCR & No se construye ni despliega & Producción intacta: Lambda lee de ECR \\
\bottomrule
\end{tabularx}
\end{table}

La disponibilidad (tabla~\ref{tab:09-dependencias}) descansa en servicios gestionados
multizona, no en configuración propia. Todo está en \code{eu-west-1}; \code{us-east-1} solo
aloja el certificado de CloudFront. Sin RPO/RTO escritos; deducidos del diseño, en \code{core} el RPO es casi continuo por PITR (35~días),
pero el RTO no se ha medido ni ensayado; \code{interactions} no tiene copia; los vectores tienen
RPO igual al último \code{just index} y RTO de minutos. El resto sale de Git y Terraform,
salvo el WAF, que se recrearía a mano. No hay \emph{runbook} de DR.

\subsection{Costes}

El coste fijo estimado es de unos 5~€/mes en recursos gestionados por Terraform (\adr{0023};
tabla~\ref{tab:06-costes}) más \mbox{$\sim$8}~\$/mes del WAF creado a mano: unos 12--13~€/mes en
total. Es una estimación; la factura de septiembre debe sustituirla. El variable lo domina Bedrock. Con Haiku~4.5 (1/5~\$ por millón de tokens de entrada/salida), un borrador de cuatro días se estima en \mbox{$\sim$0,02}~\$ y un turno de una o
dos llamadas en 0,002--0,006~\$ (\cref{sec:05-arquitectura-ia}). El coste real queda en cada
traza, pero no se ha agregado, y los tokens de embeddings no se contabilizan (TRA-232). Mil
borradores al mes costarían unos 20~\$, más que la infraestructura, y API Gateway no limita el uso
por cliente: es el principal riesgo de gasto.

\begin{table}[tb]
\centering\footnotesize
\caption{Métricas del proyecto (líneas y recuentos, 23-09-2026).}\label{tab:09-metricas}
\begin{tabularx}{\columnwidth}{@{}Yr@{}}
\toprule
Métrica & Valor \\
\midrule
Commits desde 10-03-2026 / sept. & 324 / 187 \\
PR fusionadas / sept. & 181 / 145 \\
Issues TRA citadas en commits & 107 \\
Python: backend y tools (274 arch.) & 43\,774 \\
TS/TSX: frontend (259 arch.) & 34\,623 \\
Terraform / YAML / Markdown & 2\,560 / 1\,019 / 8\,469 \\
Tests backend ejecutados & 1\,087 \\
Tests frontend ejecutados / e2e & 747 / 9 \\
ADR / \emph{runbooks} / recetas \code{just} & 23 / 7 / 41 \\
\bottomrule
\multicolumn{2}{@{}p{\columnwidth}@{}}{\scriptsize Recuentos con \code{git}, \code{gh} y
\code{wc}, sin código generado ni vendorizado; tests ejecutados con \code{just test} (todos pasan;
desglose en la \cref{tab:08-pruebas}); los 9 \emph{specs} e2e no se ejecutaron.}
\end{tabularx}
\end{table}

\subsection{Propuesta mínima}

Tres medidas casi gratuitas: AWS Budgets con aviso al 80\,\% de 30~€/mes; alarmas SNS sobre
errores y \emph{throttles} de Lambda y 5xx de API Gateway; y una ruta de salud pública sin datos
para un comprobador externo. Y dos SLO iniciales: disponibilidad mensual del API del
99,5\,\% (unas 3,6~h de presupuesto de errores al mes) y p95 del primer evento SSE del
planificador por debajo de 3~s, que ya mide la traza (\cref{sec:10-analisis-critico}).

\section{Evaluación de la solución}\label{sec:09b-evaluacion}

La \cref{tab:09b-metricas} reúne lo que hoy se mide del sistema, con su origen y su grado de
evidencia: \emph{medido} (instrumentado y reproducible), \emph{ejecutado} (pruebas que pasan o
fallan), \emph{estimado} (cálculo sobre precios de lista) o \emph{en CI} (se ejecuta en cada PR,
no para esta memoria). Las cifras de recuperación, latencia y coste proceden del spike TRA-151 y
se detallan en las \cref{tab:05-calidad,tab:05-latencia,tab:05-coste}; aquí no se repiten sus
alternativas.

\begin{table*}[tb]\centering\footnotesize
\caption{Métricas técnicas y funcionales. Fechas de 2026. Spike: \code{docs/architecture/vector-store-spike.md}; \code{ai}: \code{src/backend/services/ai\_api}; \code{cc}: \code{src/backend/tools/city\_corpus}.}
\label{tab:09b-metricas}
\setlength{\tabcolsep}{4pt}
\begin{tabularx}{\textwidth}{@{}>{\raggedright\arraybackslash}p{3.3cm}>{\raggedright\arraybackslash}p{5.6cm}>{\raggedright\arraybackslash}p{1.9cm}Y@{}}
\toprule
Métrica & Valor & Obtención & Origen y límites \\
\midrule
\multicolumn{4}{@{}l}{\textit{Técnicas}} \\
R@5 / R@10 / MRR & 0,688 / 0,812 / 0,591 (S3 Vectors) & medido, 17-09 & Spike: 32 consultas, 52 ids esperados, solo Budapest, 4 en español, sin intervalos de confianza \\
Acuerdo con coseno exacto & 0,969 del top-10 & medido, 17-09 & Spike; mide el almacén, no la relevancia \\
Latencia búsqueda (caliente) & p50 64~ms, p95 69~ms (2\,048~MB); frío 91~ms & medido, 17-09 & Spike, 128 consultas desde Lambda en \code{eu-west-1} \\
Latencia embedding Titan & $\sim$75~ms; búsqueda total p50 137~ms & medido, 17-09 & Spike; despreciable frente al LLM \\
Coste de recuperación & $\sim$2,5~USD por millón de consultas; indexar Budapest, 0,012~USD & estimado / medido & Spike; precios de \code{us-east-1} \\
Coste LLM turno / borrador & 0,002--0,006~USD / $\sim$0,02~USD (4 días) & estimado & Tokens según el código y precios de Haiku~4.5 (\cref{tab:05-llamadas}) \\
Tests backend & 37 / 199 / 597 / 254 (\code{common}, \code{core}, \code{ai}, \code{cc}), en verde & ejecutado, 23-09 & \code{just test}; LLM guionizado, DynamoDB con moto (\cref{tab:08-pruebas}) \\
Tests frontend & 747 (Vitest), en verde & ejecutado, 23-09 & \code{just test-frontend} \\
E2E & 9 \emph{specs} de Playwright & en CI & \code{src/frontend/e2e}; no ejecutados para esta memoria \\
Lint y contratos & ruff, pyright, eslint; \emph{drift} de OpenAPI y tipos & en CI & \code{.github/workflows/pr.yml}; bloquean el \emph{merge} \\
\midrule
\multicolumn{4}{@{}l}{\textit{Funcionales}} \\
Gate de \emph{readiness} & 11 umbrales; pasan Berlín, Budapest y Madrid; Bolonia no (3 hoteles con foto de 10) & ejecutado, 22-09 & \code{cc/config/readiness.py}; \code{data/<ciudad>/report.json} \\
Sesión con modelo real & 7 turnos; falla si una actividad no tiene foto o una tarjeta lleva precio & manual & \code{just planner-smoke}; Nemotron y recuperador por palabras, no Bedrock; sin registro de resultados \\
Flujos de usuario (e2e) & brief $\rightarrow$ 3 días con hotel $\rightarrow$ cambio de franja; reabrir por URL; renombrar y borrar viajes; consola admin; móvil; HTML prerenderizado & en CI & \code{planner}, \code{trips}, \code{admin*}, \code{mobile}, \code{prerender}, \code{smoke}; LLM sustituido por una sesión grabada \\
Modo demo & sesión grabada de Budapest, con aviso, si falta la ruta del planner & en CI & \code{planner.spec.ts} \\
Validación de itinerarios & avisos por día: saltos $>$3,5~km, ritmo, cerrados, precios & en cada turno & Función pura en \code{ai}; tests unitarios \\
Consola \code{/admin} & latencia p50/p95, p50 del primer evento, tokens, coste, usados/recuperados & medido (prod.) & Trazas de \adr{0024}; sin análisis publicado ni tokens de embeddings (TRA-232) \\
\bottomrule
\end{tabularx}
\end{table*}

\paragraph{Lectura.} La evidencia es sólida en lo estructural (1\,087 tests backend y 747
frontend en verde, contratos y lint que bloquean el \emph{merge}, un gate que impide indexar una
ciudad incompleta) y débil en lo que el usuario percibe. La recuperación se midió una vez, en una
ciudad y con 32 consultas: sirvió para decidir el almacén, no para certificar la calidad. Los
costes por turno son estimaciones sobre el código; la consola ya registra el coste real, pero sin
agregar por viaje ni publicar.

\paragraph{Lo que no está evaluado.} La calidad del itinerario extremo a extremo (coherencia,
distancias, interés), sin juez LLM ni humano; el comportamiento con usuarios reales, que no los
hay; la carga y las cuotas de Bedrock; la seguridad más allá del análisis estático (inyección de
\emph{prompt}, sin DAST ni \emph{pentest}); la accesibilidad, sin auditoría; y el español, con un
corpus 98,8\,\% en inglés y solo 4 consultas de cruce. El análisis y las propuestas están en la
\cref{sec:10-analisis-critico}.

\paragraph{Criterios de aceptación propuestos.} Coherentes con la \cref{tab:10-propuestas}, se
propone que bloqueen el \emph{merge} o el despliegue: R@10~$\geq$~0,75 en un banco de 100
consultas por ciudad con un 30\,\% en español; una nota media del juez LLM~$\geq$~4/5 sobre
itinerarios, validada con $\kappa\geq$~0,6 frente a 50 itinerarios anotados; p95 del primer
evento~$<$~3~s, lo que exige añadir ese percentil a la consola (hoy solo da el p50); y coste p95
por viaje~$\leq$~0,05~USD, una vez contabilizados los embeddings.

\section{Análisis crítico y propuestas de evolución}\label{sec:10-analisis-critico}

Esta sección ordena por ámbito las valoraciones repartidas por la memoria y las convierte en
propuestas priorizadas con una métrica de éxito.

\subsection{Lo que sostiene el proyecto}

Las decisiones de infraestructura se apoyaron en cifras. \adr{0014} eligió S3 Vectors por
simplicidad operativa y el \emph{spike} TRA-151, confirmatorio, construyó los dos almacenes: con
32 consultas de una ciudad (4 en español) no hubo diferencia de calidad apreciable, sin intervalos
de confianza, y el arranque en frío favoreció a S3 Vectors (\cref{sec:05-arquitectura-ia}). Las
tres versiones de la nube y la retirada de RDS se justificaron con cifras de la cuenta real
(\adr{0009}, \adr{0023}; \cref{sec:06-nube}), y la política de fotos prefiere no mostrar ninguna
antes que una equivocada (\adr{0021}, \adr{0022}). El anclaje es estructural: toda tarjeta procede
de un documento recuperado del corpus, los ids no recuperados se descartan y los precios son
niveles, nunca importes; la prosa libre del modelo (respuestas de chat, el \code{why} de 140
caracteres) sí puede mencionar lugares y no está verificada. El asistente que Expedia retiró
respondía sin inventario ni precios reales~\cite{mk:romie}; Kyrian aporta anclaje documental
(existencia, descripción, fuente), pero no disponibilidad ni precio, que exigen integrar inventario
(horizonte de 12 meses; \cref{sec:03-mercado}). Completan la base un coste fijo estimado de unos
12--13~€/mes (5~€ en Terraform más $\sim$8~\$ del WAF manual; \cref{sec:09-sre}), contratos
verificados en CI, trazas de turno (\adr{0024}), un corpus con atribución y un proceso trazable de
issue a PR y a ADR (\cref{sec:08-cicd-proceso}).

\subsection{Lo que se haría distinto}

\paragraph{Producto y experiencia.} El producto se construyó sin validación con usuarios: no hay
entrevistas, analítica de uso ni valoración por tarjeta. Lo delatan el recorte a una ciudad por
viaje (\adr{0019}) y ``Mis viajes'', añadido porque el login dejaba al usuario ante un planner
vacío (\adr{0020}). El viajero no puede corregir a la IA (TRA-146); la tarjeta no enseña la fuente
ni el fragmento recuperado, que son el diferenciador, y su contenido sale en inglés; el mapa une
paradas con líneas rectas, sin tiempos (\cref{sec:02-producto}). La accesibilidad está declarada,
no medida.

\paragraph{Inteligencia artificial.} Tres huecos pesan más que el resto:
\begin{itemize}
\item \emph{Sin evaluación extremo a extremo.} TRA-148 no tiene código: los 597 tests de
  \code{ai\_api} comprueban el orquestador con un LLM guionizado, no si un itinerario es bueno. No
  hay juez LLM~\cite{cr:zheng2023judge}, ni regresión de \emph{prompts} en CI, ni guardrails: el
  texto de Wikivoyage, editable por cualquiera, entra en el \emph{prompt} y abre un vector de
  inyección indirecta~\cite{cr:greshake2023injection}.
\item \emph{Lo medido no está en producción.} La fusión BM25+RRF subió el MRR de 0,591 a 0,651 en
  el \emph{spike}, pero las tarjetas se sirven con búsqueda densa sola (\cref{tab:05-calidad}).
\item \emph{Lo decidido no está implementado.} La política de corpus (TRA-210: $\leq$2\,000
  documentos por ciudad, cuotas y puntuación de calidad) sigue en papel: Berlín tiene 12\,240
  documentos, un 53\,\% de restaurantes sin descripción, y Bolonia 3 hoteles con foto de 118
  (\cref{tab:05-corpus}). Los \emph{gates} cuentan cantidades, no calidad: la introducción
  vandalizada de Madrid y sus guías de distrito vacías los superaron.
\end{itemize}
El corpus está en un 98,8\,\% en inglés para un público hispanohablante, y el cruce de idiomas se
midió con 4 de las 32 consultas. Hay además deuda de ingeniería: JSON pedido por \emph{prompt} con
temperatura 0,7 y una llamada de reparación en lugar de \emph{tool use}~\cite{cr:bedrocktooluse};
un único modelo, sin \emph{failover} automático; una escalera de filtros que reembebe la misma
consulta, sin caché; y tokens de embeddings sin contabilizar (TRA-232).

\paragraph{Arquitectura de software.} El doble modo de autenticación (HS256 local, RS256 con
Cognito) se ha vuelto permanente; \code{ai\_api} deja 60~min sin revocación
(\cref{sec:04-arquitectura-software}). Sin límites de uso por usuario, una cuenta comprometida o
un bucle del cliente es el mayor riesgo económico real de un sistema de $\sim$12~€/mes. La traza
se escribe en el camino crítico ($\sim$100~ms) en una tabla sin PITR, guarda texto libre 90~días
y no se borra con la cuenta. La documentación deriva: \code{code-quality-review.md} y \adr{0013}
son anteriores a DynamoDB y \code{CLAUDE.md} exige un MCP de Pencil no configurado.

\paragraph{Nube, operación y proceso.} No hay alarmas, presupuesto de AWS, comprobación externa,
\emph{staging} ni despliegue gradual; el WAF está fuera de Terraform y el estado guarda dos
secretos (\cref{sec:06-nube,sec:09-sre}): la observabilidad de negocio va por delante de la
operativa. El \emph{bus factor} es 1, los revisores son modelos con sesgos compartidos y el coste
en \emph{tokens} no se mide; las 145 PR de septiembre (\cref{tab:09-metricas}) superan la
capacidad de validar lo entregado.

\subsection{Propuestas priorizadas}

La \cref{tab:10-propuestas} agrupa las propuestas en tres horizontes. Esfuerzo e impacto son
estimaciones cualitativas; cada fila une la carencia que la justifica con su criterio de éxito.

\begin{table*}[tb]\centering\scriptsize
\caption{Propuestas priorizadas. Esf.: esfuerzo; Imp.: impacto (B bajo, M medio, A alto); Hor.:
horizonte en meses. Tras ``$\triangleright$'', la métrica de éxito.}
\label{tab:10-propuestas}
\begin{tabularx}{\textwidth}{@{}>{\raggedright\arraybackslash}p{7.2cm}lYccc@{}}
\toprule
Propuesta & Ámbito & Evidencia / métrica de éxito & Esf. & Imp. & Hor. \\
\midrule
\multicolumn{6}{@{}l}{\textbf{\color{kyrian}1 mes: proteger lo que hay}} \\
AWS Budgets al 80\,\% y alarmas CloudWatch$\to$SNS (5xx, duración, \emph{throttles}, coste de Bedrock)~\cite{cr:budgets} & Plataforma & Hoy un pico de gasto o una caída solo se ven si alguien mira (\cref{sec:09-sre}) $\triangleright$ 0 alertas sin atender en 24~h & B & A & 1 \\
\emph{Throttling} y \emph{usage plan} en API Gateway~\cite{cr:apigwthrottling}; cuota diaria por usuario; contabilizar embeddings (TRA-232) y agregar el coste por turno & Plataforma & Bedrock sin límite por usuario; coste en trazas nunca agregado $\triangleright$ coste p95 por viaje $\leq$0,05~\$ & B & A & 1 \\
Ruta de salud pública, comprobador externo y \emph{smoke test} al final de \code{deploy-backend.yml} & Plataforma & Salud tras Cognito; el primer contacto con producción es el despliegue $\triangleright$ disponibilidad $\geq$99,5\,\%; \emph{smoke} verde en cada despliegue & B & M & 1 \\
Privacidad y términos publicados, borrado de trazas al suprimir la cuenta (RGPD), aviso de IA (art.~50) y atribución CC BY-SA visible en cada tarjeta & Legal & El login enlaza políticas que no existen; trazas 90~días $\triangleright$ test de borrado en CI; 100\,\% de tarjetas con fuente & B & A & 1 \\
WAF a Terraform; secretos fuera del estado; actualizar \code{code-quality-review.md}, \adr{0013} y la regla de Pencil & Plataforma & WAF manual ($\sim$8~\$/mes); dos secretos en el estado $\triangleright$ \code{plan} sin recursos fuera de Terraform & B & M & 1 \\
\midrule
\multicolumn{6}{@{}l}{\textbf{\color{kyrian}3 meses: medir y mejorar la calidad}} \\
Banco de evaluación en CI (TRA-148): 100 consultas por ciudad, recall/MRR y juez LLM sobre itinerarios, con umbrales que bloquean el \emph{merge} & IA & 597 tests no miden calidad; eval solo de Budapest $\triangleright$ R@10 $\geq$0,75 (30\,\% en español); juez $\geq$4/5 y $\kappa\geq$0,6 frente a 50 itinerarios anotados & M & A & 3 \\
Mitigar la inyección: contenido recuperado delimitado como datos, filtro de instrucciones al construir el corpus, salida libre sin URL & IA & Wikivoyage editable; intro de Madrid vandalizada $\triangleright$ 0 éxitos en el test adversarial del banco & M & A & 3 \\
Política de corpus TRA-210 con puntuación de calidad y detector de vandalismo (edad y diff de la revisión) & IA & Berlín 6,1 veces el tope $\triangleright$ $\leq$2\,000 documentos por ciudad & M & A & 3 \\
Corpus bilingüe (traducción o embeddings evaluados en español): el segmento es hispanohablante & IA & 98,8\,\% en inglés; 4 consultas en español $\triangleright$ R@10 en español a 0,05 del inglés & A & A & 3 \\
Híbrido BM25+RRF para tarjetas; embeber la consulta una vez y cachear; \emph{reranker}~\cite{cr:bedrockrerank} solo si el banco lo justifica & IA & MRR 0,591$\to$0,651 (\cref{tab:05-calidad}) $\triangleright$ MRR $\geq$0,65 sin perder R@10 & M & M & 3 \\
\emph{Tool use}~\cite{cr:bedrocktooluse}; temperatura 0--0,3 para JSON; Sonnet para el ranking solo si el banco compensa su coste (2--3$\times$) & IA & JSON por \emph{prompt} y reparación $\triangleright$ 0 llamadas de reparación & M & M & 3 \\
Pruebas de carga con k6~\cite{fz:k6} o Locust~\cite{fz:locust} contra Compose con LLM simulado y prueba corta real contra Bedrock & Plataforma & Sin pruebas de carga; cuotas de Bedrock sin medir $\triangleright$ usuarios concurrentes antes de \emph{throttling}; p95 primer evento $<$3~s & M & A & 3 \\
Edición manual y guardado (TRA-146); fuente y fragmento en la tarjeta, en el idioma del usuario; valoración por tarjeta ligada a la traza & UX & El viajero no puede corregir a la IA ni ver la fuente $\triangleright$ $\geq$4/5 de valoración media & A & A & 3 \\
Auditoría axe~\cite{cr:axe} y Lighthouse en CI & UX & WCAG~2.2 declarado, no medido $\triangleright$ Lighthouse a11y $\geq$95 & B & M & 3 \\
\emph{Staging} por \emph{workspace} de Terraform y alias de Lambda con \emph{canary}~\cite{cr:codedeploy} & Plataforma & Un solo entorno $\triangleright$ 0 \code{apply} sin \code{plan} probado en \emph{staging} & M & M & 3 \\
\midrule
\multicolumn{6}{@{}l}{\textbf{\color{kyrian}12 meses: ampliar producto y negocio}} \\
Piloto B2B con una DMO en marca blanca, sobre \code{/add-city} & Negocio & LTV/CAC B2B de 6--12 frente a $\approx$0,5 en B2C (\cref{sec:03-mercado}) $\triangleright$ un contrato firmado & A & A & 12 \\
Viajes multiciudad; rutas y tiempos entre paradas con OSRM~\cite{cr:osrm} o Valhalla~\cite{cr:valhalla}; exportar a PDF/ICS y compartir & UX & \adr{0019}; líneas rectas; sin salida del producto $\triangleright$ tiempo entre paradas en el 100\,\% de los días & A & A & 12 \\
Agente con herramientas para datos vivos (horarios, precios, disponibilidad): conserva la procedencia y solo muestra un importe con fuente en vivo citada, la excepción a ``nunca un número'' de \adr{0018} & IA & Sin inventario ni precios (el fallo de Romie) $\triangleright$ 100\,\% de importes con fuente & A & A & 12 \\
Afiliación medida (Booking, GetYourGuide) y plan comercial de Open-Meteo & Negocio & 0,15--0,76~€ por viaje, hipotético; Open-Meteo no comercial~\cite{mk:openmeteo} $\triangleright$ ingreso por viaje medido & M & M & 12 \\
Revisión humana muestral de PR y coste por ola; retirar o portar \code{infra/gcp}; segunda región solo si un SLO lo exige & Proceso & \emph{Bus factor} 1; \emph{tokens} sin medir $\triangleright$ 10\,\% de PR revisadas por una persona & B & M & 12 \\
\bottomrule
\end{tabularx}
\end{table*}

\subsection{Hoja de ruta}

El orden sigue un criterio explícito. El primer mes acota el riesgo económico, operativo y legal:
presupuesto, alarmas, límites de uso y privacidad cuestan horas y cierran las vías por las que el
proyecto puede hacerse caro o ilegal de golpe. Los tres meses siguientes miden la calidad de la IA
antes de mejorarla: sin banco de evaluación no se puede saber si la búsqueda híbrida, un
\emph{reranker} o un modelo mayor mejoran algo. Solo entonces tiene sentido ampliar el producto y
probar el negocio con un piloto B2B, que además daría los primeros usuarios reales.

\section{Conclusiones y trabajo futuro}\label{sec:11-conclusiones}

Kyrian World funciona en producción con un coste fijo estimado de $\approx$12~€/mes
($\approx$5~€/mes en Terraform más el WAF creado a mano; sin factura real todavía) y anclaje
estructural de sus tarjetas; su calidad de itinerario, su aceptación por usuarios y su viabilidad
económica siguen sin medir. Un equipo de cuatro personas, con agentes de IA y un proceso escrito, lo
llevó a producción en seis meses para cuatro ciudades. El anclaje tiene un alcance preciso: toda
tarjeta procede de un documento recuperado y los ids no recuperados se descartan, pero la prosa
libre del modelo (respuestas de chat, el \code{why} de 140 caracteres) puede mencionar lugares no
verificados. El \emph{spike} de vectores confirmó S3 Vectors, elegido antes por simplicidad
operativa (\adr{0014}): sin diferencia de calidad apreciable en 32 consultas de una ciudad.

No hay usuarios reales, así que no se sabe si el viajero valora la verificabilidad lo bastante
como para preferirla a un chat generalista gratuito, y la viabilidad económica descansa en
hipótesis (\cref{sec:03-mercado}). El trabajo futuro, priorizado, está en la
\cref{sec:10-analisis-critico}; antes de ponerlo delante de usuarios de pago habría que hacer
cuatro cosas:
\begin{enumerate}[nosep]
\item Acotar el riesgo: límites de uso por usuario, presupuesto y alarmas.
\item Medir la calidad: un banco de evaluación extremo a extremo en CI que bloquee regresiones.
\item Cumplir en privacidad: publicar la política y los términos y borrar las trazas al suprimir
  la cuenta.
\item Validar con usuarios: valoración por tarjeta y un piloto con una DMO que aporte uso real y
  un primer ingreso.
\end{enumerate}
El vídeo demo opcional que contempla el guion no se incluye en esta entrega.


{\scriptsize\bibliographystyle{unsrtnat}
\bibliography{referencias}}
\appendix
\section{Glosario y referencia de decisiones}\label{sec:12-anexo}
{\footnotesize
\paragraph{Glosario.}
\begin{description}[font=\normalfont\bfseries,leftmargin=0pt,itemsep=0pt]
\item[ADR] \emph{Architecture Decision Record}: documento numerado con contexto, decisión y consecuencias; no se borra, se sustituye.
\item[Base de datos vectorial] almacén que indexa \emph{embeddings} y devuelve los más cercanos a una consulta (aquí, S3 Vectors).
\item[Brief] especificación escrita antes del código: resultado, decisiones cerradas, ficheros, pruebas y aceptación.
\item[Cold start] arranque en frío: latencia extra de la primera invocación de una Lambda recién creada.
\item[DMO] \emph{Destination Management Organisation}: organismo público o mixto que promociona un destino.
\item[Embedding] vector numérico que representa el significado de un texto; textos parecidos quedan cerca.
\item[Gate de readiness] comprobación automática que decide si el corpus de una ciudad está completo para publicarse.
\item[Grounding] anclaje: toda tarjeta o lugar propuesto por el modelo debe existir en el corpus recuperado.
\item[IaC] infraestructura como código: los recursos de AWS se declaran en Terraform y se versionan en Git.
\item[MRR] \emph{Mean Reciprocal Rank}: media de $1/p$, con $p$ la posición del primer resultado correcto.
\item[OIDC] OpenID Connect: federación con la que GitHub Actions obtiene credenciales temporales de AWS sin claves.
\item[RAG] generación aumentada por recuperación: el modelo responde con documentos recuperados en el \emph{prompt}.
\item[Recall@$k$] fracción de consultas con al menos un documento esperado entre los $k$ primeros.
\item[RRF / BM25] BM25: ranking léxico por términos; RRF: fusión de rankings por suma de $1/(k+p)$.
\item[SSE] \emph{Server-Sent Events}: flujo HTTP unidireccional de eventos del servidor al navegador.
\item[Tabla única] diseño de DynamoDB con todas las entidades en una tabla, distinguidas por claves compuestas.
\item[Turno / traza] turno: una petición del usuario y su respuesta; traza: su registro con pasos, tokens, coste y latencia.
\item[Worktree] copia de trabajo adicional de Git sobre otra rama, aislada de la principal.
\end{description}

\par\medskip\noindent\begin{minipage}{\columnwidth}\textbf{Decisiones de arquitectura} (\code{docs/architecture/adr/}; la 0012 no existe). A: aceptada; S: sustituida; P: en parte.
\par\smallskip
\begin{tabularx}{\columnwidth}{@{}lYl@{}}
\toprule
ADR & Decisión & Estado \\
\midrule
0001 & Backend en \code{core\_api} y \code{ai\_api} & A \\
0002 & JWT sin estado; reenvío del token & S por 0009 \\
0003 & Frontend con dos URL base & A \\
0004 & \code{src/}, \code{infra/}, \code{docs/} & A \\
0005 & \code{Trip} raíz de agregado & A (S-P 0019) \\
0006 & Modelo de vista en el frontend & A (S-P 0019) \\
0007 & AWS; SSO local, OIDC en CI & A \\
0008 & AWS v2: CloudFront, API GW, Fargate & S por 0009 \\
0009 & AWS v3: Lambda, Cognito, 30~€/mes & A (datos: 0023) \\
0010 & Zona DNS y certificado en Terraform & A \\
0011 & Viajes reales, carga en cliente & A (S-P 0019) \\
0013 & Conversaciones en \code{core\_api} & A \\
0014 & Vectores en S3 Vectors & A \\
0015 & Planificador SSE v2 sin estado & A \\
0016 & Mapa MapLibre + OpenFreeMap & A \\
0017 & Intro y foto de ciudad en el manifiesto & A \\
0018 & Respuestas de chat con tarjetas & A \\
0019 & Los viajes viven en el planificador & A (enm. 0020) \\
0020 & Inicio con sesión: página de viajes & A \\
0021 & Vista previa del sitio del local & A \\
0022 & Foto de hotel resuelta al construir & A \\
0023 & DynamoDB sustituye a RDS & A (enm. 0024) \\
0024 & Trazas de turno y acceso admin & A \\
\bottomrule
\end{tabularx}\end{minipage}

\paragraph{Issues citadas.}
TRA-28 despliegue en GitHub Pages (\code{.nojekyll});
TRA-133 imágenes a ECR con \code{crane};
TRA-146 borrador del planificador como viaje;
TRA-148 banco de evaluación de Budapest;
TRA-151 \emph{spike} Qdrant frente a S3 Vectors;
TRA-180 WebGL en el devcontainer;
TRA-187 planificador en el \emph{viewport} móvil;
TRA-195 dependencias, tres mayores retenidas;
TRA-210 corpus curado con cuotas y fotos;
TRA-219 retirada de RDS en dos \code{apply};
TRA-204 Miami a través del pipeline; TRA-232 tokens de \emph{embeddings} en la traza.
\par}

\end{document}
